\begin{dcomment}{exposition}{p139}{The replacement must be taken over the new base}
Grothendieck now has a particularly economical candidate for direct image: a category over $S'$ is regarded as a category over $S$. The difficulty is to show that this operation has the expected universal property after localization. Its prospective right adjoint is not arbitrary strict pullback. It is pullback after a replacement which remembers arrows in the base. The next seven manuscript pages make that distinction precise enough to construct an adjunction, including its two triangle identities.

Write $Q_S=\Phi_S$ and $i_X:X\longrightarrow Q_SX$. For a category $p:X\longrightarrow S$, an object of $Q_SX$ is a triple $(s,x,a:s\longrightarrow p(x))$. A morphism from $(s,x,a)$ to $(t,y,b)$ consists of $v:s\longrightarrow t$ and $u:x\longrightarrow y$ satisfying $p(u)a=bv$. The projection to $s$ is the split fibration associated with the contravariant functor $s\longmapsto s\backslash X$. The map $i_X$ sends $x$ to $(p(x),x,\id_{p(x)})$. This recalls the comma replacement discussed at manuscript p.~20 (``Which comma construction repairs which variance''), but here its role is more demanding: it supplies the unit of a derived adjunction.

For $f:S'\longrightarrow S$, denote by $L$ the functor induced by composition with $f$, and by $R$ derived inverse image. If $[X]$ denotes $\hot_{S,W}(X)$, their formulas are
\[
 L[X']=[X'/S],\qquad R[X]=[Q_SX\times_SS'].
\]
The replacement in the second formula is over $S$. Even when $X'\longrightarrow S'$ is already a fibration, its composite with an arbitrary $f$ need not be one. Replacing over $S'$ would therefore solve the wrong problem. This is the reason for the otherwise rather insistent occurrences of $X'/S$ in the argument. The occurrence of $Q_S(X'/S')$ in the square on manuscript page 139 is a slip already present in the edition; its intended meaning is $Q_S(X'/S)$.

The unit can be written before localization. Its component is the unique $S'$-functor
\[
 a_{X'}:X'\longrightarrow Q_S(X'/S)\times_SS'
\]
whose first projection is $i_{X'/S}$ and whose second projection is the structure map of $X'$. Explicitly, it sends $x'$ to the triple with base object $f(p'(x'))$ and identity arrow, together with $p'(x')\in S'$. The construction is natural in $X'$. Passing through the localization therefore gives $\alpha:\id\longrightarrow RL$.

The counit is slightly different. If $X\longrightarrow S$ is smooth, the projection $X\times_SS'\longrightarrow X$ represents $\beta_{[X]}:LR[X]\longrightarrow[X]$. For an arbitrary representative, the formula is instead
\[
 \beta_{[X]}=[i_X]^{-1}
 [\pr_1:Q_SX\times_SS'\longrightarrow Q_SX].
\]
Both arrows in this fraction lie over $S$, and the arrow inverted is one already known to belong to $W_S^{\mathrm d}$. The forgetful functor $Q_SX\longrightarrow X$ on underlying categories does not generally lie over $S$: the two structure maps at $(s,x,a)$ have values $s$ and $p(x)$. The fraction is thus doing necessary work.

The formulas also explain why naturality survives the change of representatives. For an arrow $u:X\longrightarrow Y$ over $S$, the equalities $Q_S(u)i_X=i_Yu$ and $Q_S(u)\pr_1=\pr_1(Q_S(u)\times_SS')$ give the naturality square for $\beta$ after the $i$'s have been inverted. The construction uses the localization universal property, rather than a choice of inverses in the original category. This is precisely the amount of structure needed for the adjunction proof that follows.
\end{dcomment}

\begin{dcomment}{examples}{p142}{A pullback can destroy the replacement equivalence}
Grothendieck's first attempt to construct the counit asks whether the map $i_X$ can be inverted after an arbitrary pullback. The walking arrow shows exactly where that attempt breaks. We work with ordinary weak homotopy equivalences, $W=W_\infty$, so that an empty category cannot be equivalent to a point.

Let $S=[1]=(0\longrightarrow1)$ and let $X=\{1\}\longrightarrow S$. An object of $\Phi_SX$ is an arrow of $S$ with target $1$. There are two such objects, $0\longrightarrow1$ and $\id_1$, and one nonidentity morphism between them. The projection identifies $\Phi_SX$ with $S$ over itself. Consequently
\[
 i_X:\{1\}\longrightarrow S
\]
is right-fine: for $s=0$ and $s=1$, the functor $s\backslash X\longrightarrow s\backslash S$ is a weak equivalence between contractible categories.

Now pull back along $j:\{0\}\longrightarrow S$. The resulting map is
\[
 X\times_S\{0\}=\varnothing
 \longrightarrow
 \Phi_SX\times_S\{0\}=\{0\}.
\]
It is not in $W$, nor is it a right-fine equivalence over the point. Regarded again as a map over $S$, it still fails the right-fine test at $0$. Thus the proposed inverse of $\pi_X=i_X\times_S\id_{S'}$ need not exist in either localization relevant to the attempted argument.

The calculation also distinguishes the two inverse images without any abstract machinery. Strict pullback of $X$ to $0$ is empty, whereas derived pullback of $[X]$ to $0$ is a point. The missing object records the arrow $0\longrightarrow1$. A comma fibre retains that arrow; a strict fibre does not.

There is no defect in the replacement itself. The map $i_X$ becomes an isomorphism over $S$, exactly as intended. The invalid move was to apply a functor which had not been shown to preserve that class of equivalences on all of $\Cat/S$. On manuscript page 143 Grothendieck repairs the argument himself: he first projects to $\Phi_SX$, and then inverts $i_X$ in the localization over $S$. No pullback of an arbitrary denominator is required.
\end{dcomment}

\begin{dcomment}{exposition}{p146}{The two triangles reduce to ordinary naturality}
The adjunction proof is complete, although the second triangle occupies several diagrams. Its purpose is to turn the chosen representative formulas into a genuine universal property. We give both checks together, keeping the distinction between identities in a category of categories and identities after localization.

For $X'$ over $S'$, put $Z=\Phi_S(X'/S)$ and let $a:X'\longrightarrow Z\times_SS'$ be the unit representative. The first triangle is immediate from the defining equality $\pr_1a=i_{X'/S}$:
\[
 \beta_{L[X']}L(\alpha_{[X']})
 =[i_{X'/S}]^{-1}[\pr_1][a]=\id_{L[X']}.
\]
This calculation uses a denominator over $S$, where its invertibility was established. In particular, it does not revive the failed pullback argument of page 142.

For the second triangle we may represent an arbitrary object of $\HOT_W(S)$ by a smooth $X\longrightarrow S$. Put $X'=X\times_SS'$ and $p=\pr_1:X'\longrightarrow X$. The map
\[
 r_X=i_X\times_S\id_{S'}:
 X'\longrightarrow\Phi_SX\times_SS'
\]
is now right-fine over $S'$. Here the smoothness of both $X$ and $\Phi_SX$ is essential: right-fine equivalences between these good representatives are preserved by arbitrary base change. This is a property of the representatives, not a property of arbitrary arrows in $\Cat/S$.

Under these identifications, the description of the derived counit says
\[
 [r_X]\circ R(\beta_{[X]})
 =[\Phi_S(p)\times_S\id_{S'}].
\]
The unit $a_{X'}:X'\longrightarrow\Phi_SX'\times_SS'$ therefore satisfies
\[
 [r_X]\circ R(\beta_{[X]})\alpha_{R[X]}
 =[\Phi_S(p)\times_S\id_{S'}][a_{X'}].
\]
The right side is $[r_X]$. To see this before localization, compare its two projections. Projection to $S'$ is the structure map of $X'$. Projection to $\Phi_SX$ gives $\Phi_S(p)i_{X'}$, which equals $i_Xp$ by naturality of $i$. These are exactly the projections of $r_X$. Since $[r_X]$ is invertible, cancellation proves
\[
 R(\beta_{[X]})\alpha_{R[X]}=\id_{R[X]}.
\]

The target of $r_X$ is $\Phi_S(X)\times_SS'$. In the displayed line immediately preceding the first triangle on manuscript page 144 the edition prints $\Phi_S(X')\times_SS'$ instead. The adjacent diagram and the expression $i_X\times_S\id_{S'}$ determine the intended target unambiguously. This is a typographical slip, not a gap in the naturality argument.

We can now use $L\dashv R$ without further representative choices. Its Hom bijection sends $h:L\xi'\longrightarrow\eta$ to $R(h)\alpha_{\xi'}$, and its inverse sends $k:\xi'\longrightarrow R\eta$ to $\beta_\eta L(k)$. The two equations just proved say exactly that these transformations are inverse. They also make the construction independent of an accidental presentation of a morphism by a fraction. The result is stronger than the assertion that direct and inverse image have compatible formulas on objects.
\end{dcomment}

\begin{dcomment}{context}{p147}{Adjunction is the organizing operation}
Grothendieck's insistence on constructing both adjunction arrows belongs to the original motivation of derivators. In \emph{Pursuing Stacks}, \S69, printed pp.~1--2, he describes the diagram categories as a way to recover structure lost by a single derived category, and identifies the two adjoints to restriction as the operations replacing direct and inverse limits \citep{Grothendieck1983}. The present calculation supplies one of those adjoints in the categorical model of homotopy types.

Its immediate use is already visible. Once the unit and counit exist, a natural transformation between diagrams of indexing categories can be transported across adjunctions. This produces the comparison maps for Kan extensions, rather than merely plausible isomorphisms between their values. The commentary on Chapter~I, manuscript p.~31 (``The base-change map as a mate'') explained that formal operation; manuscript pages 138--146 have now constructed the particular adjunction to which it will be applied.

There is a further gain. An adjunction may pass to a reflective subcategory even when its two functors are not constructed there in the same manner. Direct image descends through a localization, while inverse image restricts to local objects. The next argument separates these two mechanisms. That separation will become decisive when the chapter later distinguishes $\HOT_W^{\mathrm{lc}}$ from $\HOT_W^{\mathrm{lc}!}$.
\end{dcomment}

\begin{dcomment}{exposition}{p150}{Passing an adjunction through two reflections}
Grothendieck wants the locally constant adjunction without repeating the categorical calculations. The abstract argument is correct, and it is useful to state its hypotheses independently of any identification of local homotopy theories. Suppose
\[
 \rho:\mathcal A\rightleftarrows\mathcal A_0:i,
 \qquad
 \rho':\mathcal A'\rightleftarrows\mathcal A'_0:i'
\]
are adjunctions with $i,i'$ fully faithful, and $L:\mathcal A'\rightleftarrows\mathcal A:R$ is an adjunction. Let $\Sigma$ and $\Sigma'$ be the classes inverted by $\rho$ and $\rho'$. Assume that $L(\Sigma')\subseteq\Sigma$ and that $R$ carries the essential image of $i$ into that of $i'$. We then have functors $\bar L$ and $\bar R$, with natural isomorphisms
\[
 \bar L\rho'\simeq\rho L,
 \qquad i'\bar R\simeq Ri.
\]
The desired adjunction follows from the natural chain
\[
\begin{aligned}
 \Hom_{\mathcal A_0}(\bar L\rho' x',y)
 &\simeq\Hom_{\mathcal A_0}(\rho Lx',y)\\
 &\simeq\Hom_{\mathcal A}(Lx',iy)\\
 &\simeq\Hom_{\mathcal A'}(x',Riy)\\
 &\simeq\Hom_{\mathcal A'}(x',i'\bar Ry)\\
 &\simeq\Hom_{\mathcal A'_0}(\rho'x',\bar Ry).
\end{aligned}
\]
Every object of $\mathcal A'_0$ is isomorphic to one of the form $\rho'x'$. The resulting bijection also descends naturally in that object. Precomposition with a localization is fully faithful on functor categories: a naturality square for an inverted arrow is obtained from the original square by inversion, and then naturality holds for every composite of ordinary arrows and formal inverses. This supplies the step which cannot be replaced by essential surjectivity alone. If one wishes to avoid conventions about the notation $\mathcal A_0^{\wedge}$ in the text, the relevant object is simply the bifunctor on $(\mathcal A'_0)^{\mathrm{op}}\times\mathcal A_0$ appearing in this display.

The two compatibility assumptions actually describe the same condition. This standard observation clarifies why Grothendieck calls them dual. An object $z\in\mathcal A$ belongs to the reflective subcategory if and only if
\[
 \Hom(b,z)\longrightarrow\Hom(a,z)
\]
is bijective for every $a\longrightarrow b$ in $\Sigma$. One direction follows from adjunction. For the converse, apply the condition to the unit $z\longrightarrow i\rho z$. It gives a retraction $i\rho z\longrightarrow z$. Since $i\rho z$ is already local, the same orthogonality shows that the other composite is the identity. Thus the unit is invertible.

If $L$ preserves localization equivalences and $y$ is local, the adjunction identifies
\[
 \Hom(b',Ry)\longrightarrow\Hom(a',Ry)
\]
with $\Hom(Lb',y)\longrightarrow\Hom(La',y)$. This is a bijection for every $a'\longrightarrow b'\in\Sigma'$, so $Ry$ is local. Conversely, if $R$ preserves local objects, the same argument shows that $\rho L(a'\longrightarrow b')$ induces bijections into every object of $\mathcal A_0$. Yoneda then makes it invertible. In particular, one need not verify two unrelated conditions in applications. The basic facts about reflections and fully faithful adjoints are recorded in \citet[Lemma~4.5.11 and Definition~4.5.12]{Riehl2016}; the displayed Hom calculations give the particular descent argument needed here.

This formal lemma should be carried unchanged into the later revision of the chapter. What changes there is its application. A coarse localization of categories over $S$ does not automatically identify with the full subcategory of locally constant objects. The latter comparison requires the additional hypotheses examined on manuscript pages 171--204. Theorem~2 therefore gives the fine adjunction without that difficulty, and gives the locally constant adjunction whenever the asserted reflections and their compatibility have been established. A correct formal argument does not supply a missing identification of its input categories.
\end{dcomment}

\begin{dcomment}{exposition}{p151}{Deriving the product in both variables}
The product appears because Grothendieck wants to understand the endofunctor $f_!f^*$. Before identifying it with multiplication by a relative homotopy type, he has to decide which fibre products survive localization. Proposition~4 makes the safe construction: take both representatives smooth over the base.

For smooth $X,Y\longrightarrow S$, their fibre product is again smooth. If $u:X\longrightarrow X_1$ is a right-fine equivalence between smooth objects, pulling it back along an arbitrary $Y\longrightarrow S$ gives a right-fine equivalence over $Y$. Forgetting from $Y$ to $S$ preserves that class. Hence
\[
 X\times_SY\longrightarrow X_1\times_SY
\]
is right-fine over $S$. The same assertion holds in the other variable when the objects varied are smooth. Localizing separately in the two variables therefore produces a bifunctor on $\HOT_W(S)$. The uniqueness asserted in the proposition is uniqueness with its comparison isomorphisms, rather than literal equality of all choices of localized categories.

For arbitrary representatives its value is
\[
 [X]\times_S[Y]
 \simeq[\Phi_SX\times_S\Phi_SY].
\]
If $X$ is smooth, $\Phi_SX$ may be replaced by $X$ in this expression. That fact does not yet allow $\Phi_SY$ to be replaced by $Y$. The two replacements concern different maps. The first is an equivalence between smooth objects, to which the base-change statement applies. The second would require arbitrary pullback along $X\longrightarrow S$ to preserve all right-fine equivalences. Smoothness of $X$ alone does not say this.

Associativity, symmetry and the unit are transported with their coherence maps. Their diagrams commute already for strict fibre products of categories; localization preserves those equalities. The unit is the class of $S\longrightarrow S$. Equivalently, the diagonal--product adjunction and the terminal-object adjunction descend, so the multiplication is the categorical product in the localized category. In the usual homotopy theory its diagram interpretation is simply the pointwise product of space-valued presheaves.

The boxed formula preceding the proposition must consequently be read provisionally. A formula
\[
 f_!f^*(\xi)\simeq[S'/S]\times_S\xi
\]
for arbitrary $f$ would identify a construction deriving only the varying input with one deriving both factors. Grothendieck notices the discrepancy almost immediately. The following counterexample is therefore part of the proof's development, and the corrected corollary is its actual conclusion. Treating the crossed-out assertion as a theorem would obscure the main mathematical distinction these pages establish.
\end{dcomment}

\begin{dcomment}{examples}{p154}{The false corollary on every fibre}
Grothendieck chooses the smallest base on which the two product constructions can differ. We expand his example because it also identifies the endofunctor which replaces the false formula. Let $S=[1]=(0\longrightarrow1)$, let $X=\{0\}\longrightarrow S$, and let $Y=\{1\}\longrightarrow S$.

The inclusion of $0$ is a discrete fibration: an arrow into $0$ can only be its identity. Thus $X$ is smooth. Its comma fibres are
\[
 0\backslash X=e,\qquad1\backslash X=\varnothing.
\]
For $Y$ both comma fibres are points, and $\Phi_SY\cong S$. The relative homotopy type $[Y]$ is therefore the terminal object of $\HOT(S)$. It follows that
\[
 [X]\times_S[Y]\simeq[X].
\]
On the other hand, $X\times_SY=\varnothing$, so
\[
 [X\times_SY]=[\varnothing].
\]
Evaluation at $0$ distinguishes the two objects: it gives a point on the first side and the empty homotopy type on the second. One smooth factor has not made the strict fibre product into the derived product.

There is a useful computation for an arbitrary coefficient. Write a space-valued presheaf on $[1]$ as a map $b:B\longrightarrow A$, with values $F(1)=B$ and $F(0)=A$. It may be represented by a smooth category over $S$, or calculated in the usual space-valued diagram model. For the inclusion $j_1:\{1\}\longrightarrow S$, restriction gives $B$, and left Kan extension returns the constant diagram on $B$. Thus
\[
 j_{1!}j_1^*F=(B\xrightarrow{\id}B),
 \qquad [Y]\times_SF=F=(B\xrightarrow{b}A).
\]
The counit has components $\id_B$ at $1$ and $b$ at $0$. The candidate product formula is an equivalence precisely when $b$ is a weak equivalence. For $F=[X]$, this is the map $\varnothing\longrightarrow e$.

For comparison, the inclusion $j_0:\{0\}\longrightarrow S$ gives
\[
 j_{0!}j_0^*F=(\varnothing\longrightarrow A)
 \simeq[X]\times_SF.
\]
The difference between the endpoints has a categorical explanation. The object $0$ defines a sieve, so left Kan extension of presheaves from it is extension by the empty object. The object $1$ instead has an incoming arrow from outside its image, and that arrow transports $B$ into the value at $0$.

This also proves a sharper warning contained in the corrected corollary. The map $Y\longrightarrow S$ is a right-fine equivalence, so $[Y]\simeq[S]$. Nevertheless the two fixed-factor endofunctors are different: multiplication by the category $S/S$ is the identity, whereas $\multhot_{Y/S}$ sends $B\longrightarrow A$ to the constant diagram on $B$. Thus $\multhot_{Y/S}$ depends on the chosen category $Y$ over $S$, not only on its class in $\HOT(S)$.
\end{dcomment}

\begin{dcomment}{development}{p155}{A precise boundary for the product formula on finite posets}
The walking-arrow example suggests asking which inclusions permit the formula for every coefficient. In a finite poset the answer can be read directly from the order. The following elementary development extends Grothendieck's example to a complete criterion in that setting. We use ordinary homotopy types throughout, so $W=W_\infty$.

\textbf{Proposition (full subposets and the universal product formula).}\ours
Let $S$ be a finite poset, regarded as a category, and let $u:U\hookrightarrow S$ be the inclusion of a full subposet. In space-valued presheaves put $P_U=u_!1$, where $1$ is the terminal presheaf on $U$. The canonical map
\[
 \theta_F:u_!u^*F\longrightarrow P_U\times F
\]
is an equivalence for every $F:S^{\mathrm{op}}\longrightarrow\mathsf{Spaces}$ if and only if $U$ is a lower ideal: whenever $s\le t$ and $t\in U$, one has $s\in U$.

Here $u_!$ means homotopy left Kan extension along $u^{\mathrm{op}}$. The two components defining $\theta_F$ are $u_!u^*F\longrightarrow u_!1$, induced by the terminal map, and the counit $u_!u^*F\longrightarrow F$. This specifies the comparison whose invertibility is claimed.

\emph{Proof.} For $s\in S$, write
\[
 U_{\ge s}=\{t\in U\mid s\le t\}.
\]
The pointwise formula gives
\[
 (u_!G)(s)\simeq
 \operatorname*{hocolim}_{(U_{\ge s})^{\mathrm{op}}}G(t),
 \qquad
 P_U(s)\simeq B(U_{\ge s}),
\]
where $B$ denotes the nerve's homotopy type. Replacing a category by its opposite does not change this homotopy type.

Suppose first that $U$ is a lower ideal. If $s\notin U$, the category $U_{\ge s}$ is empty, so both $(u_!u^*F)(s)$ and $(P_U\times F)(s)$ are empty. If $s\in U$, the poset $U_{\ge s}$ has initial object $s$, and its opposite has terminal object $s$. The homotopy colimit is therefore $F(s)$; the same argument with constant value $e$ gives $P_U(s)\simeq e$. Under these identifications $\theta_F(s)$ is the identity of $F(s)$. This proves the sufficiency, including the comparison map, at every object of $S$.

For necessity, suppose that $s\notin U$ but $s\le t_0$ for some $t_0\in U$. Take the representable discrete presheaf
\[
 F=h_s,\qquad h_s(t)=\Hom_S(t,s).
\]
For every $t\in U_{\ge s}$, one has $s\le t$ and $t\ne s$. Antisymmetry then excludes $t\le s$, so $h_s(t)=\varnothing$. The diagram computing $(u_!u^*h_s)(s)$ is constantly empty, and its homotopy colimit is empty. By contrast,
\[
 (P_U\times h_s)(s)\simeq B(U_{\ge s})\times e
\]
is nonempty, since $t_0$ is an object of $U_{\ge s}$. Hence $\theta_{h_s}(s)$ is not an equivalence. This contradicts the asserted universal formula and proves necessity.\hfill$\square$

The test uses only a representable presheaf. Consequently, in this finite-poset setting the universal formula can already be tested on the collection of representables $h_s$. For a point inclusion $\{t\}\hookrightarrow S$, it holds for every coefficient exactly when $t$ is minimal. It holds for every point inclusion exactly when $S$ is discrete.

There is also a distinction from absolute weak equivalence. In the chain $0<1<2$, take $U=\{1,2\}$. Both $U$ and $S$ have contractible nerves, and $P_U\simeq1$: each nonempty upper intersection $U_{\ge s}$ has a least element. Nevertheless $u^*h_0$ is empty, so $u_!u^*h_0$ is empty while $P_U\times h_0\simeq h_0$ is not. Equivalence of the total homotopy types of the indexing categories does not remove the variation of an arbitrary coefficient diagram.

The positive sieve calculation is standard Kan-extension calculus; compare \citet[Proposition~1.23, with the indexing categories replaced by their opposites]{Groth2013}. We checked the pointwise and sieve formulations in the cited derivator literature and searched for the corresponding finite-poset converse under projection formula and Frobenius reciprocity. We did not locate this exact elementary criterion there. The proof is included as a continuation of the example, without a claim of priority or of a new general projection-formula theorem.
\end{dcomment}

\begin{dcomment}{exposition}{p156}{Fixed-factor multiplication is a direct-image composite}
The corrected corollary does more than impose an extra hypothesis on the false formula. It isolates an operation which remains available for every category $x:X\longrightarrow S$. Its varying argument is derived, while the category $X$ is held fixed:
\[
 M_X([Y])=[X\times_S\Phi_SY].
\]
The base-change assertion used in Proposition~4 makes this independent of the chosen smooth representative of $[Y]$. The example just calculated shows why it is not generally independent of the representative of $[X]$.

This operation has a direct adjunction-theoretic description:
\[
 M_X\simeq x_!x^*.
\]
Indeed, for smooth $Y$ the inverse image $x^*[Y]$ is represented over $X$ by $Y\times_SX$. Direct image along $x$ regards the same total category as lying over $S$. Thus both sides give $[Y\times_SX]$, naturally in $Y$. Since smooth objects present the localization, the identification holds on all of $\HOT_W(S)$. Corollary~2 is the special case $X=S'$.

There is a canonical comparison with the derived product. The map $i_X:X\longrightarrow\Phi_SX$ induces, for smooth $Y$,
\[
 [X\times_SY]\longrightarrow[\Phi_SX\times_SY]
 \simeq[X]\times_S[Y].
\]
The right side derives both inputs; the left side has only derived the varying input. If $X$ is smooth, the comparison is invertible, because it is the pullback of a right-fine equivalence between smooth representatives. Without that hypothesis the comparison can fail, as it does for $X=\{1\}$ over the walking arrow. The later discussion of homotopical fibrations and properness will supply further sufficient hypotheses. Nothing in the present corollary establishes those extensions in advance.

The identity $M_X\simeq x_!x^*$ also supplies structure which is easy to miss if one sees only the fibre-product formula. It is the underlying endofunctor of a comonad. Its counit is the projection map already used in the adjunction proof, and its comultiplication is
\[
 x_!x^*\xrightarrow{x_!\alpha x^*}x_!x^*x_!x^*.
\]
The comonad laws follow from the triangle identities. When $X$ is smooth, this is multiplication by $[X]$ with the diagonal coalgebra structure coming from the cartesian product. This interpretation explains the order of the argument: the adjunction is established first, and the formula for its composite is then refined according to the geometry of the fixed factor.

Thus the three cases listed on manuscript pages 155--156 have different information content. With two smooth factors, one obtains the product of two objects of the homotopy category. With only the varying factor smooth, one obtains the value of $M_X$. For two arbitrary strict representatives, their strict fibre product need not be determined by either localized class. These are statements about the operations available from particular representatives, rather than about different possible definitions of the same product.
\end{dcomment}

\begin{dcomment}{exposition}{p157}{The comma square and its actual comparison map}
Grothendieck now verifies the pointwise formula for direct image. This is where the construction becomes a derivator calculation: evaluation at an object of the base is expressed by a homotopy colimit over a category of arrows. The comma replacement makes the formula almost tautological, but the comparison map and the variance deserve to be stated.

Fix $s\in S$ and write $C=s\backslash S'$. Its objects are pairs $(s',a:s\longrightarrow f(s'))$, and a morphism from $(s',a)$ to $(t',b)$ is $v:s'\longrightarrow t'$ such that $f(v)a=b$. Let $q:C\longrightarrow S'$ forget $a$, let $r:C\longrightarrow e$, and let $p:e\longrightarrow S$ select $s$. The square drawn in the text is not strictly commutative. It carries the canonical natural transformation
\[
 \tau:pr\Longrightarrow fq,
 \qquad \tau_{(s',a)}=a.
\]
The two cartesian squares in the larger diagram involving $s\backslash S$ are, by contrast, strict pullback squares. These two kinds of square should not be identified.

With the presheaf convention $\HOT(S)=\operatorname{Ho}(\mathsf{Spaces}^{S^{\mathrm{op}}})$, $\tau$ gives $q^*f^*\longrightarrow r^*p^*$. The standard left-adjoint mate is
\[
 r_!q^*\xrightarrow{r_!q^*\alpha}
 r_!q^*f^*f_!\longrightarrow
 r_!r^*p^*f_!\xrightarrow{\beta}
 p^*f_!.
\]
It points in the reverse direction to the arrow displayed on manuscript page 157. Once exactness is proved, that displayed arrow can be read as its inverse. The direction is fixed by the units, counits and structural $2$-cell; it is not a further hypothesis on $W$. For the general mate construction, see \citet[Proposition~4.3.8]{Riehl2016}.

Let $X'\longrightarrow S'$ represent $\xi'$. The right side $p^*f_!\xi'$ is represented by $s\backslash X'$. On the left, $q$ is a discrete opfibration: an arrow $v:s'\longrightarrow t'$ has the unique lift taking $(s',a)$ to $(t',f(v)a)$. The pullback rule for these co-isomorphisms permits us to compute $q^*\xi'$ by strict pullback even for arbitrary $X'$. That pullback has objects $(a:s\longrightarrow f(s'),x')$ with $p'(x')=s'$, and is canonically isomorphic to $s\backslash X'$. Applying $r_!$ then forgets its map to $C$, so
\[
 r_!q^*\xi'\simeq\hot_W(s\backslash X')
 \simeq p^*f_!\xi'.
\]

These identifications also calculate the mate. An object represented by $(a,x')$ is carried by the unit into the comma replacement with its identity arrow; the structural transformation then supplies $a$; the counit forgets the redundant intermediate object. The resulting object is precisely $(a,x')$ in $s\backslash X'$, and the same description sends a morphism to itself. Thus the mate is the identity under these categorical identifications. This proves the canonical comparison is invertible, rather than merely exhibiting an unspecified isomorphism between its source and target.

In modern diagram notation the formula reads
\[
 (f_!F)(s)\simeq
 \operatorname*{hocolim}_{(s\backslash f)^{\mathrm{op}}}F(s').
\]
The opposite is forced by the presheaf convention. Modern presentations using covariant diagrams write the corresponding over-category formula instead. The pointwise axiom and its comma-square formulation appear in \citet[Definition~1.5 and Proposition~1.26]{Groth2013}; the calculation here supplies them concretely for the categorical model.
\end{dcomment}

\begin{dcomment}{exposition}{p158}{Constant coefficients detect the total homotopy type}
Having obtained direct images, Grothendieck can express an equivalence of indexing categories by its effect on homology with constant coefficients. This assertion concerns the total homotopy type of the base. It should be kept separate from equivalence of all categories of diagrams over the base.

Let $p:S\longrightarrow e$ and represent $\xi\in\Hot_W$ by a category $T$. Since every category is smooth over a point, the preceding product computation gives
\[
 p_!p^*(\xi)\simeq\hot_W(S)\times\xi.
\]
For $f:S'\longrightarrow S$, with $p'=pf$, the induced comparison is therefore
\[
 \hot_W(S')\times\xi\longrightarrow\hot_W(S)\times\xi.
\]
The subscript $S$ on the multiplication signs in the edition's display on page 158 is a slip: both objects belong to $\Hot_W=\HOT_W(e)$, so this is the product over $e$.

If $f\in W$, its image in $\Hot_W$ is invertible; applying the product functor proves the desired comparison is invertible for every $\xi$. Conversely, if every such comparison is invertible, take $\xi=e$, the product unit. Then $\hot_W(f)$ itself is invertible. Thus the precise condition detected is membership in the strong saturation $\overline W$, the class of arrows inverted by $\Cat\longrightarrow\Hot_W$. When $W$ is already strongly saturated this is simply $f\in W$. The later restatement on manuscript page 204 makes this distinction explicit.

The walking arrow again illustrates the scope. The inclusion $\{1\}\longrightarrow[1]$ is a weak equivalence of nerves. It therefore gives the same homology of constant coefficients as $[1]\longrightarrow e$. Restriction of arbitrary presheaves along that inclusion, however, remembers only $B$ from a diagram $B\longrightarrow A$. It cannot be an equivalence of all diagram categories: the diagrams $\varnothing\longrightarrow e$ and $\varnothing\longrightarrow\varnothing$ have the same restriction but are distinguished at $0$.

The stronger statement about \emph{locally constant} diagrams is different. On $[1]$, local constancy requires $B\longrightarrow A$ to be an equivalence, so evaluation at either endpoint loses no information up to coherent equivalence. This small example explains why the next page strengthens the conclusion only on $\HOT_W^{\mathrm{lc}}$. It also explains why the axiom controlling the relation between coarse equivalence and fibrewise equivalence will enter that proof. The elementary product argument for constant coefficients does not itself establish that stronger assertion.
\end{dcomment}

\begin{dcomment}{status}{p159}{What Theorem 3 establishes at this point}
Grothendieck has proved the left-adjoint and pointwise formulas, and has checked homological invariance for constant coefficients. These are substantial conclusions. The annotations to Theorem~3 also record two separate qualifications which must remain part of its statement.

First, the claim of Der~2 is withdrawn in footnote~76. A family of evaluations is conservative if a morphism in the localized diagram category is invertible whenever all its evaluations are invertible. Defining weak equivalences objectwise does not, by itself, prove this statement for arbitrary morphisms of a localization. One still has to control how such morphisms are represented, or compare the localization with a model in which objectwise detection can be proved. The supplement on manuscript page 164 belongs to that unfinished part of the argument. It would be inaccurate either to count the parenthetical ``pour mémoire'' as a proof or to infer that conservativity is false in ordinary homotopy theory.

For the standard model the question is settled. \citet[Theorem~6.11, including its proof of Der~2]{Cisinski2003} constructs homotopy right Kan extensions for model categories with small limits and proves the corresponding derivator axioms; his Proposition~6.13 constructs the space-valued external Hom. For combinatorial model categories, \citet[Proposition~1.30]{Groth2013} gives a strong derivator. These results justify the usual space-valued diagram interpretation used in the examples above.

For every basic localizer, Der~2 now follows from \citet[Propositions~3.1.30 and~4.2.7]{Cisinski2006}: conservativity for presheaf localizers transfers through a diagramwise equivalence using $\Delta$ (4.1.29). Thus the revised relative W$(1,2,3)$ suffice, without accessibility or W(4). This later repair preserves the historical withdrawal; it makes no assertion for arbitrary weak basic localizers.

Second, the strengthened assertion that a $W$-equivalence of bases induces an equivalence of locally constant categories uses additional control of fibres. The annotation names the earlier axiom L6, and the following pages examine exactly where it is needed. For ordinary weak equivalences, homotopy invariance of locally constant model-category diagrams is proved by \citet[Theorem~1.3]{Cisinski2009LocallyConstant}.

The final paragraph about constructing $f_*$ together with internal Hom is therefore a programme at this place in the text. The adjunction proved above supplies $f_!$, but it supplies no right adjoint to $f^*$ by formal symmetry: taking opposites of indexing categories is not the same as changing a left adjoint into a right adjoint in the same coefficient theory. The later model-category constructions give both operations in their stated settings. The localizer-by-localizer analysis in the remainder of this chapter remains necessary to determine what follows from Grothendieck's chosen axioms alone.
\end{dcomment}

\begin{dcomment}{exposition}{p161}{When a change of base becomes an equivalence}
Grothendieck now wants to know how much of the base a locally constant coefficient can remember. The adjunction is already available; the task is to determine when its two comparison maps become invertible. This is a stronger assertion than the agreement of cohomology with constant coefficients discussed just before it. A locally constant family has transport, and changing the base must preserve the whole family if restriction is to be an equivalence.

Write $f:S'\to S$, and retain the hypotheses under which the locally constant adjunction was constructed. For a perfect representative $X\to S$, the counit is the projection
\[
 \varepsilon_X:X\times_S S'\longrightarrow X.
\]
The notation $\beta_0$ used for this map agrees with the preceding adjunction. To prove that $\varepsilon_X$ is a $W$-equivalence from the hypothesis $f\in W$, one must know that this particular change of base preserves $W$. This is the role of L6. The fact that a product with a fixed category preserves $W$ is insufficient: $X\to S$ need not be a product projection.

The unit has a different presentation because an arbitrary $X'\to S'$ need not remain a good representative after composition with $f$. Put $P=C_S(X')$. The structural map $P\to S$ is perfect, and the natural map $j:X'\to P$ is a total $W$-equivalence. The map
\[
 a:X'\longrightarrow P\times_S S',\qquad
 \operatorname{pr}_1a=j,
\]
is the unit on the chosen representative. If $f\in W$, L6 gives $\operatorname{pr}_1\in W$, and two-out-of-three gives $a\in W$. Thus both the unit and the counit are inverted. The triangle identities then make $f_!^0$ and $f_0^*$ mutually quasi-inverse equivalences. The label $\beta_{0,X'}$ in the edition is a slip: this second comparison is $\alpha_{0,X'}$. The same pages twice place $\hot^{\mathrm{lc}}$ in an ambient category without the superscript $\mathrm{lc}$; the locally constant category is intended.

The footnote about necessity also has a precise reading. Suppose that the adjunction on locally constant objects is an equivalence for every $f\in W$. Its counit then shows that $X\times_S S'\to X$ is inverted by the relevant localization for every perfect $X\to S$. Applying the total-space functor shows that it is inverted in $\Hot_W$. If $W$ is strongly saturated, it belongs to $W$. Consequently the asserted invariance forces the perfect-pullback stability used in the proof. Without strong saturation the same argument gives membership in the strong saturation of $W$, which explains why the manuscript cautiously writes ``$\pm$ n\'ecessaire''.

This argument will remain useful after the axioms are reorganized on manuscript p.~171. It identifies exactly the property needed at this point, without treating an earlier choice of notation for that property as permanent. It also distinguishes two operations which the succeeding pages repeatedly separate: replacing the total category by a better object over $S$, and testing whether that replacement is still an equivalence on each fibre.
\end{dcomment}

\begin{dcomment}{exposition}{p163}{The canonical map in the base-change theorem}
The purpose of Theorem~4 is to turn the concrete model for $f_!$ into a calculational rule. A pullback square of categories gives two ways of transporting a coefficient from $X$ to $S'$. An isomorphism between their values would not by itself suffice: one needs the particular comparison determined by the adjunction, since subsequent compositions of squares use its compatibility with pasting.

For the square in the text, the comparison is the composite
\[
 f'_!q^*\xrightarrow{\ f'_!q^*\eta_f\ }
 f'_!q^*f^*f_!
 \cong f'_!f'^*p^*f_!
 \xrightarrow{\ \varepsilon_{f'}p^*f_!\ }p^*f_!.
\]
Here $q^*f^*\cong f'^*p^*$ is induced by the equality $fq=pf'$. This formula fixes the orientation. In the smooth case choose $Z\to X$ smooth. Then $Z\to S$ is smooth, and all inverse images in this expression can be computed literally. The unit introduces the relevant pullback and the counit projects from it; their composite is the canonical identification between the two copies of $Z\times_S S'$. Thus the calculation on p.~162 proves that the canonical map is invertible, and not merely that its source and target happen to be isomorphic.

The proper case uses a different good feature. If literal pullback along $p$ carries $W_S^{\mathrm d}$ into $W_{S'}^{\mathrm d}$, then for every $Z\to S$ the comparison
\[
 Z\times_S S'\longrightarrow \Phi_S(Z)\times_S S'
\]
is right-fine. Hence literal pullback computes $p^*$ even for a representative which is not smooth over $S$. To finish, one may still choose $Z\to X$ smooth in order to compute $q^*$; this costs nothing, because every coefficient has such a representative. Both sides of base change are again represented by $Z\times_S S'$, and the displayed mate gives its identity after the canonical comparisons. The direct argument on pp.~168--169 proves precisely the preservation statement needed here.

A small example shows why smoothness concerns transport rather than constancy. Let $S=(0\to1)$ and take the presheaf of sets $F$ with
\[
 F(0)=\{a,b\},\qquad F(1)=\{c\},\qquad c\longmapsto a.
\]
Its category of elements $X$ has the arrow $a\to c$ and the isolated object $b$, with the evident map $f:X\to S$. This is a discrete Grothendieck fibration, hence smooth. For the terminal coefficient on $X$, the object $f_!1_X$ is represented by $X\to S$. Pullback to $0$ is the two-point category, whereas pullback to $1$ is the one-point category. Both values agree with base change along the corresponding point inclusion. The transition between the fibres is not a weak equivalence, and no local constancy has been used.

The terminology is inherited from the base-change theorems of algebraic geometry, but the hypotheses here are conditions on categories of lifts. The precise categorical characterization is recorded by \citet[Theorem~1.2]{Maltsiniotis2005Asphericite}; its proper version follows by opposites. In the modern derivator formalism, \citet[Propositions~1.24 and~1.26]{Groth2013} treats the corresponding exact squares for Grothendieck fibrations, opfibrations and comma categories. These references clarify the later language without supplying the missing Der~2 argument in the manuscript's first attempted proof.
\end{dcomment}

\begin{dcomment}{exposition}{p165}{What constant coefficients detect}
Grothendieck interrupts the base-change argument to identify the absolute equivalences detected by $\HOT_W$. This is a useful move: it prevents the phrase ``$\HOT_W$-equivalence'' from concealing a new, unexplained class of maps. The proof uses the adjunction just established and the already available product formula for a constant coefficient.

Let $p_S:S\to e$ and put $\eta=\hot_W(S)$. For $a\in\Hot_W$, write $a_S=p_S^*a$. The representing-category description gives a natural isomorphism
\[
 p_{S!}(a_S)\cong\eta\times a.
\]
Indeed, if $a=\hot_W(T)$, then $a_S$ is represented by $S\times T\to S$ and its left direct image is represented by $S\times T$. This also identifies the morphism induced by $f:S'\to S$: it is $\hot_W(f)\times\id_a$.

Adjunction therefore turns the defining map for a $\HOT_W$-equivalence into
\[
 \Hom_{\Hot_W}(\eta\times a,b)
 \longrightarrow
 \Hom_{\Hot_W}(\eta'\times a,b),
 \qquad \eta'=\hot_W(S').
\]
The base-dependent subscripts in the edition's displayed Hom formula are slips; both Hom sets here are in $\Hot_W$. If $\hot_W(f)$ is an isomorphism, every displayed map is a bijection. Conversely, assume these maps are bijections for all $a,b$, and take $a$ to be the terminal object $1$. The product identifications $\eta\times1\cong\eta$ and $\eta'\times1\cong\eta'$ reduce the assertion to bijectivity of
\[
 \Hom_{\Hot_W}(\eta,b)\longrightarrow\Hom_{\Hot_W}(\eta',b)
 \quad\hbox{for every }b.
\]
Yoneda gives invertibility of $\hot_W(f)$. More explicitly, taking $b=\eta'$ supplies a left inverse of $\hot_W(f)$, and injectivity for $b=\eta$ forces it to be a right inverse. This proves Proposition~5.

Strong saturation enters only when the result is translated back from the localized category to the original class $W$. The proposition without saturation says that the detected maps form $\overline W$, the class inverted by $\Cat\to\Hot_W$. If $W=\overline W$, this is exactly $W$. The proof does not require the detection of all isomorphisms in $\HOT_W(S)$ by point evaluations. In particular, it must not be counted as a proof of Der~2. That latter statement concerns arbitrary arrows of a localized category, including zigzags, whereas Proposition~5 identifies the image of a given functor of indexing categories.

The source itself recognizes the distinction on p.~166. The intended appeal to the abstract proper-base-change theorem still needs Der~2, so the proof by that route is incomplete. Nothing is wrong with Proposition~5, and nothing forces Theorem~4 to wait for a global repair of the derivator axioms: the direct proof on pp.~168--169 works before that repair. Keeping this order of dependence is the reason to retain both the interrupted argument and the author's later replacement of it.
\end{dcomment}

\begin{dcomment}{exposition}{p167:Comme ce foncteur transforme $W$ en $W$, on gagne.}{A single coefficient already tests right-fine maps}
The lemma on pp.~166--167 is meant to extract a statement about actual functors from a base-change isomorphism. Its family of constant coefficients is natural from the cohomological viewpoint. Once the pointwise formula for $f_!$ is known, however, the terminal coefficient already detects the maps under consideration. This observation shortens the proof and makes its independence from Der~2 explicit.

Assume $W$ is strongly saturated, and let $u:X'\to X$ be a morphism over $S$. Then
\[
 u\in W_S^{\mathrm d}
 \quad\Longleftrightarrow\quad
 \hot_{S,W}(u)\text{ is an isomorphism}.
\]
The forward implication is the definition of the localization. For the converse, apply the evaluation functor $s^*$ to the asserted isomorphism. The pointwise formula identifies its value with
\[
 \hot_W(s\backslash X')\longrightarrow\hot_W(s\backslash X).
\]
Strong saturation makes $s\backslash u$ a member of $W$ for every $s$, which is exactly the definition of $W_S^{\mathrm d}$. This argument applies evaluation to an isomorphism already known to exist; it does not infer a global isomorphism from a family of pointwise isomorphisms. Consequently it uses no unproved conservativity axiom.

Now let $p:X\to S$ and $p':X'\to S$. The image of the terminal coefficient under left direct image is $p_!1_X=\hot_{S,W}(X)$, and the map induced by $u$ is $\hot_{S,W}(u)$. Thus the preceding equivalence is precisely the manuscript's test with $\xi=1$. For a general constant coefficient represented by $T$, the induced map is represented by $u\times\id_T$. Since
\[
 s\backslash(X\times T)\cong(s\backslash X)\times T,
\]
product stability of $W$ proves the implication from $u\in W_S^{\mathrm d}$ to the test for every $\xi$. The larger family of tests is therefore equivalent to the single one, rather than an additional hypothesis.

This gives Corollary~1 with its quantifiers visible. Suppose that base change along $p:S'\to S$ is valid for left direct image along every $X\to S$. If $u:X_1\to X_2$ belongs to $W_S^{\mathrm d}$, its map of pushed-forward terminal coefficients is invertible. Pulling that map back along $p$ preserves invertibility. Naturality of base change identifies the resulting map with the left direct image of the terminal-coefficient map induced by
\[
 u':X_1\times_S S'\longrightarrow X_2\times_S S'.
\]
The single-coefficient test over $S'$ implies $u'\in W_{S'}^{\mathrm d}$. Strong saturation is used at this last conversion to membership in $W$; without it the argument establishes the corresponding statement for the class detected in the localization.

This is a useful instance of the distinction between a criterion for a representative and a criterion for an arbitrary localized morphism. The former is enough for the present base-change theorem. The latter is the broader problem of Der~2, which remains separate.
\end{dcomment}

\begin{dcomment}{exposition}{p169}{Proper base change without a pointwise conservativity axiom}
The short proof on pp.~168--169 succeeds because it turns a relative question into an absolute one over a suitable comma category. Properness supplies a comparison of bases which remains a $W$-equivalence after every pullback. Two-out-of-three can then transport the desired equivalence across that comparison.

We first isolate the elementary part. Let $v:T'\to T$ be a universal $W$-equivalence, meaning that for every $A\to T$ the projection $A\times_TT'\to A$ belongs to $W$. If $h:A\to B$ is a map over $T$, the commutative square
\[
\begin{tikzcd}
 A\times_TT' \arrow[r,"h'"]\arrow[d,"v_A"']&
 B\times_TT'\arrow[d,"v_B"]\\
 A\arrow[r,"h"']&B
\end{tikzcd}
\]
has both vertical arrows in $W$. If $h\in W$, then $hv_A\in W$ and $v_Bh'=hv_A$ forces $h'\in W$. The reverse implication is the same calculation in the other direction. Literal pullback along a universal equivalence therefore preserves and reflects $W$. No localization, no adjunction and no pointwise detection is involved in this lemma.

Take now $g:S'\to S$ proper, or more generally $W$-proper, and fix $s'\in S'$ with $s=g(s')$. Properness gives a universal $W$-equivalence
\[
 v:s'\backslash S'\longrightarrow s\backslash S.
\]
Its fibre over an arrow $s\to t$ is the category of its lifts starting at $s'$. The defining asphericity conditions for properness, together with their stability after base change, give the stated universality. This is the geometric content of the proof; it is not a consequence of the existence of initial objects in the two comma categories alone. Two contractible categories can be connected by a functor which fails to be a universal equivalence.

For $u:X\to Y$ over $S$, form $u_s:s\backslash X\to s\backslash Y$, a map over $s\backslash S$. If $u\in W_S^{\mathrm d}$, then $u_s\in W$. Pullback along $v$ identifies $u_s$ with
\[
 s'\backslash(X\times_SS')\longrightarrow
 s'\backslash(Y\times_SS').
\]
The elementary lemma makes this map a member of $W$. Since $s'$ was arbitrary, literal pullback along $g$ preserves right-fine equivalences. This completes the proper case of Theorem~4 by the model calculation explained above.

For a discrete opfibration, the role of the lift categories is especially transparent. An arrow $g(s')\to t$ has a unique lift with prescribed source $s'$. Hence the induced functor on undercategories is an isomorphism; the preceding comparison needs no homotopy argument at all. For an ordinary Grothendieck opfibration, a chosen cocartesian lift is initial in the corresponding category of lifts, so its nerve is contractible. The general proper case replaces this initial object by the weaker asphericity condition required by $W$.

The manuscript calls the earlier obstruction a ``blocage''. Mathematically, the change is a change of test objects: the universal comparison over $s\backslash S$ carries exactly the information needed for right-fine equivalence. The abandoned four-page route is not a counterexample to proper base change, nor does it expose a defect in the construction of $f_!$.
\end{dcomment}

\begin{dcomment}{exposition}{p171:Ces trois axiomes sont suffisants}{Axioms for the localization and axioms for its fibres}
Grothendieck's new beginning separates the construction of a homotopy theory from the stronger rules which make its total spaces determine its fibres. Much of the preceding work survives under three localizer axioms. The extra fibration axiom enters at identifiable comparison maps, and the next ten pages are an effort to locate those maps rather than to replace the entire construction.

The localization axiom W(2) is a relative Theorem~A statement. For a triangle $X\to Y\to S$, it says that membership of $X/s\to Y/s$ in $W$ for every $s$ implies membership of $X\to Y$ in $W$. Together with the terminal-object axiom and the saturation rules, it gives the familiar local-to-global mechanism. The weakened W(2~bis) keeps only the case in which all the relevant comma categories are $W$-aspherical. The distinction matters when the minimum class $W_\omega$ is introduced: it is defined using this weaker axiom, not by simply repeating the complete list W(1,2,3).

There is a second distinction in the word ``fibration''. The ordinary category $\Fib S$ consists of Grothendieck fibrations over $S$. The category $\Fib_W S$ in the revised diagram consists of weak $W$-fibrations and is defined by a homotopical condition. One equivalent test is this: for every categorically contractible $T\to S$ and every $t\in T$, the inclusion
\[
 X_t\longrightarrow X\times_ST
\]
belongs to $W$. This is the formulation of Chapter~XII, Theorem~3, manuscript pp.~51--52, to which Grothendieck refers here \citep{Grothendieck1990XIIa}. The word ``contractible'' in this test cannot be replaced silently by ``$W$-aspherical''. Doing that asks for a stronger condition, whose relation to the weak condition is precisely part of the fibration-axiom problem.

For an ordinary Grothendieck fibration, the weak condition has a concrete interpretation: each transport functor $u^*:X_{s'}\to X_s$ must belong to $W$. Thus an ordinary fibration can be smooth without being weakly $W$-fibred, as in the two-fibre example above. Conversely, a weak $W$-fibration need not possess chosen cartesian lifts. The comma replacement $\Phi$ is what supplies an ordinary split fibration representing its homotopical data.

W(4b) compares two tests for a morphism between such weakly fibred objects. The global test asks whether $X\to Y$ lies in $W$ as a functor of total categories. The fibre test asks whether each $X_s\to Y_s$ lies in $W$. The implication from the fibre test to the global test already follows from W(2), after comparing fibres with comma fibres. W(4b) asserts the converse. It is therefore a cancellation rule for total spaces over a fixed base, not merely a restatement of two-out-of-three.

The first version of W(4b$'$) restricts that rule to $W_\omega$-perfect objects. This does prove equality of the two tests inside that smaller category. It does not yet compare an arbitrary weakly fibred $X$ with its perfect replacement $C(X)$, because $X$ need not belong to the smaller category. The footnote retracting the claimed sufficiency points to exactly this missing scope in the available argument. The retraction does not establish that the implication is false. With the later minimality and filtered-colimit results explained below, the alternating replacement on pp.~192--195 supplies the missing fine unit; the original restricted condition does then imply W(4b). The commentary at p.~195 gives the transfer argument. Later a different condition receives the same label W(4b$'$); the two formulations must be kept separate throughout the audit.
\end{dcomment}

\begin{dcomment}{later}{p172:Le rôle que j'ai fait jouer}{The proposed theory below ordinary weak equivalence}
Grothendieck is trying to remove an externally imposed reference to classical homotopy. His question about a localizer strictly smaller than $W_\infty$ concerns the strength of the axioms themselves: perhaps their formal consequences already contain more classical homotopy theory than he has yet extracted.

For the modern basic-localizer axioms, this question is settled. \citet[Definition~1.1.3 and Theorem~2.2.11]{Cisinski2004Minimal} proves that every basic localizer contains the nerve weak equivalences $W_\infty$. Thus no strictly smaller basic localizer satisfies that complete package. The weaker list defining $W_\omega$ requires a further theorem. \citet[Theorem~1.10 and Remark~1.11]{Cisinski2009LocallyConstant} proves the corresponding minimality among weak basic localizers, using the absolute Theorem~A condition. It follows that $W_\omega=W_\infty$ as well. To check the saturation hypothesis, if $ri=\id$ and $ir\in W$, the manuscript's condition on the two composites $ri,ir$ gives $i,r\in W$; this is the retract clause in the weak-basic-localizer formulation. W(2~bis) and W(3) supply its other two clauses. The two axiom lists remain different even though their minimum classes now have the same description.

This result changes the status of the proposed example, but it leaves the reason for the revision intact. Removing an explicit hypothesis $W_\infty\subseteq W$ can still clarify which finite arguments use only the localizer axioms and which use a separate fibre-detection condition. Indeed, a class larger than $W_\infty$ can satisfy the basic axioms while failing W(4b). The connected-components localizer used below gives an elementary instance. The ordinary theory is minimal in the order of classes of equivalences; further localization is entirely possible, and it can forget information that a fibrewise test retains.

The same minimum is obtained after adjoining W(4b): $W'_\omega=W_\infty$. Indeed, Quillen's Theorem~B identifies the fibres of a weak $W_\infty$-fibration with its homotopy fibres \citep[Theorem~2.3.2]{Cisinski2004Minimal}. A total weak equivalence over a fixed base therefore induces weak equivalences on the fibres. Thus $W_\infty$ satisfies W(4b) and is an admissible candidate for $W'_\omega$; basic minimality gives the reverse inclusion. Its strong saturation is the same class. The homotopy-pullback description also supplies the strongest base-change property intended in footnote~84.

This identifies the minimum within the standing W(1,2,3) framework. The inclusion $W'_\omega\subseteq W$ does not, by itself, imply that $W$ inherits every closure property appearing in the defining list: the notion of weak $W$-fibration changes when $W$ changes. The fibre-cancellation axiom is not simply an upward-closed property of subsets of $\mathrm{Fl}(\Cat)$. This explains why Grothendieck treats the suggested containment and the assertion W(4b) as different conditions in the paragraph itself.

By \citet[Corollary~4.2.22]{Cisinski2006}, W(6~bis) is automatic for basic localizers; the theorem assumes relative W(2).
\end{dcomment}

\begin{dcomment}{exposition}{p174:Les relations entre [?] localiseurs}{Five tests and the role of good objects over the base}
The diagram of localizers is the bookkeeping on which the revised theory depends. Grothendieck is distinguishing the information carried by a strict fibre, a comma fibre, and a total category. To read its equalities, every test must refer to the same map $u:X\to Y$ over the same base $S$.

The five conditions are
\[
\begin{aligned}
 u\in W_S&\iff u\in W\text{ on total categories},\\
 u\in W_S^{\mathrm f}&\iff u_s:X_s\to Y_s\in W\quad(s\in S),\\
 u\in W_S^{\mathrm d}&\iff s\backslash X\to s\backslash Y\in W\quad(s\in S),\\
 u\in W_S^{\mathrm g}&\iff X/s\to Y/s\in W\quad(s\in S),\\
 u\in W_S^{\mathrm{univ}}&\iff
 X\times_ST\to Y\times_ST\in W\quad(T\to S).
\end{aligned}
\]
In the last line universality is relative to $S$. It is not the assertion that $u:X\to Y$ remains a $W$-equivalence after every map into its total target $Y$. The difference is already visible when $S=e$.

If $X$ and $Y$ are $W$-smooth over $S$, the inclusions of strict fibres into under-comma fibres give a commutative square
\[
\begin{tikzcd}
 X_s\arrow[r,"u_s"]\arrow[d]&Y_s\arrow[d]\\
 s\backslash X\arrow[r,"s\backslash u"']&s\backslash Y
\end{tikzcd}
\]
whose vertical arrows lie in $W$. Two-out-of-three identifies the fibre and right-fine tests. Smoothness survives change of base. Thus the fibre test, applied after every $T\to S$, gives a right-fine map over $T$ and hence a total $W$-equivalence by W(2). This proves $W_S^{\mathrm f}\subseteq W_S^{\mathrm{univ}}$ on smooth objects. The reverse inclusion follows by taking $T=\{s\}$. Finally, taking $T=S/s$ gives $W_S^{\mathrm{univ}}\subseteq W_S^{\mathrm g}$, and the relative localizer axiom gives $W_S^{\mathrm g}\subseteq W_S$. These are exactly the first chain of implications in the text. Opposites give the proper chain.

For weak $W$-fibrations, the defining universal comparisons apply to both the initial and the terminal comma categories. Hence both vertical comparisons, for $s\backslash X$ and for $X/s$, are in $W$, and remain so after changing the base. The same proof identifies all four fine or universal tests. W(4b) adds the reverse implication from the total test. Its statement is equivalent to equality of all five classes on $\Fib_W S$ for every $S$.

A two-object example explains the restriction to good objects. Let $S=(0\to1)$, let $X$ be the discrete category on $0,1$ mapped identically on objects to $S$, and let $Y=S$ with its identity structural map. The functor $u:X\to Y$ is the identity on objects. Every strict fibre map is the identity of a point, so $u\in W_S^{\mathrm f}$ for ordinary weak equivalences. But $X$ has two connected components and $Y$ has one, so $u\notin W_S$. At $0$, the under-comma category $0\backslash X$ consists of two isolated objects, whereas $0\backslash Y$ is an interval. The right-fine test detects exactly the missing arrow. This is why the text avoids using the strict-fibre localizer on all of $\Cat/S$.

The phrase that the inclusion is strict ``si on ne fait pas'' W(4b) should be read as a statement about what the basic axioms do not force. For a fixed $W$ and a fixed $S$, absence of an additional assumption is not a proof that the classes differ. The exact universal assertion is the equivalence with W(4b); the example below supplies a localizer and a base where the difference is real.
\end{dcomment}

\begin{dcomment}{exposition}{p176:induit par l'inclusion}{Why four localizations appear}
Grothendieck's large diagram organizes choices of representatives. The smaller categories make constructions explicit; the larger ones make stability properties easier to state. The object represented should not depend on a convenient replacement, but that independence must be proved for the particular class of arrows one inverts.

There are two first choices. Localizing all of $\Cat/S$ by the under-comma test gives $\HOT_W(S)$; using the over-comma test gives $\HOT_W^o(S)$. The replacements $\Phi$ and $\Psi$ make these two categories accessible through smooth and proper models respectively. Their symmetry comes from taking opposite categories, not from an identification of the two comma tests on arbitrary objects over $S$.

The weakly fibred objects make a second pair of choices possible. On $\Fib_W S$, all the fine tests agree, so one may form
\[
 \mathcal L_S=\Fib_W S[(W_S^{\mathrm f})^{-1}].
\]
This retains the homotopy type of each fibre together with its transport. One may instead invert every map of weakly fibred objects whose total functor belongs to $W$:
\[
 \mathcal G_S=\Fib_W S[W_S^{-1}].
\]
The notation in the text is $\mathcal L_S=\HOT_W^{\mathrm{lc}}(S)$ and $\mathcal G_S=\HOT_W^{\mathrm{lc}!}(S)$. The exclamation mark records that this second construction may identify more arrows. It does not add an extra adjunction or an additional kind of pointwise constancy.

Since $W_S^{\mathrm f}\subseteq W_S$ on weakly fibred objects, the identity functor induces $\lambda:\mathcal L_S\to\mathcal G_S$. This is a localization by its universal property. Namely, a functor out of $\mathcal L_S$ factors through $\lambda$ precisely when it inverts the images of all total $W$-equivalences between weakly fibred objects. No full faithfulness of $\lambda$ follows from this description.

The replacement $C$ explains why the coarse category also has the presentation
\[
 \mathcal G_S\simeq(\Cat/S)[W_S^{-1}].
\]
It sends every object into $\Parf_\omega S\subseteq\Fib_W S$, and its natural unit is a total $W$-equivalence. Therefore both composites of the inclusion and replacement become naturally isomorphic to the identity after coarse localization. The same proof applies to a full intermediate category containing the image of $C$. For fine localization, however, the proof requires a fibrewise unit; the coarse unit cannot be used in its place.

For $\Phi_0$ and $\Psi_0$ the necessary finer comparison is available on $\Fib_W S$. This is why the two wings of the upper diagram still describe $\mathcal L_S$, whereas the perfect vertex is temporarily excluded from that assertion. The pattern is not an arbitrary list of exceptions. Each exception records a precise natural transformation whose known class is too coarse for the desired localization.

The diamond on p.~178 consequently says $\rho i=\lambda=\sigma j$, up to the canonical identifications of presentations. Calling $i$ and $j$ inclusions expresses the intended role of locally constant objects inside the two directional theories. Their full faithfulness still requires an argument. The reflection formulas established after W(4b) is imposed provide such an argument; the shape of the diamond alone does not.
\end{dcomment}

\begin{dcomment}{development}{p178}{An explicit failure of faithfulness for the coarse comparison}
Grothendieck's four-category diagram raises a question about arrows, not only about which objects become isomorphic. The coarse localization can erase distinct automorphisms of a locally constant family. We make that loss explicit for the connected-components localizer. The covering obstruction discussed at manuscript p.~43 already distinguished a total equivalence from a fibre equivalence; the calculation here follows the comparison functor $\lambda$ itself and determines the localization of a whole class of examples.

Let $W_0$ be the class of functors inducing bijections on connected components. It satisfies the basic localizer axioms. Saturation and the terminal-object condition follow directly from the functor $\pi_0:\Cat\to\mathrm{Set}$. For the relative axiom, given $X\to S$, there is a natural bijection
\[
 \pi_0X\cong\mathop{\mathrm{colim}}_{s\in S}\pi_0(X/s).
\]
To see this without a homotopy-colimit theorem, send the class of an object $(x,p(x)\to s)$ to the component of $x$. Every component occurs by using $(x,\id_{p(x)})$. The relations induced by morphisms in a slice, and by arrows of $S$, preserve that component. Conversely, every object $(x,p(x)\to s)$ is identified in the colimit with $(x,\id_{p(x)})$, and every arrow $x\to x'$ gives the corresponding identification through the slice over $p(x')$. These observations give inverse maps. A pointwise bijection of the slice-component diagrams therefore induces a bijection on total components, proving W(2).

\textbf{Proposition (orbit collapse).}\ours\quad
Let $G$ be a nontrivial group and $S=BG$. For $W=W_0$, the functor
\[
 \lambda:\HOT_{W_0}^{\mathrm{lc}}(BG)
 \longrightarrow\HOT_{W_0}^{\mathrm{lc}!}(BG)
\]
is not faithful. On the full category of discrete fibrations over $BG$ associated with right $G$-sets, its fine and coarse source presentations are respectively the category of right $G$-sets and the category of sets; the induced functor is the orbit functor.

\emph{Proof.} For a right $G$-set $E$, let $\int E\to BG$ be its category of elements. Its objects are elements of $E$, and a morphism over $g\in G$ with target $e'$ has source $e'g$. Unique lifting, and the invertibility of $g$, make the projection both a discrete fibration and a discrete opfibration. The categories of prescribed lifts are singletons, so these objects are $W_\omega$-perfect and in particular weakly $W_0$-fibred. Functors over $BG$ between them correspond exactly to equivariant maps of $G$-sets.

The strict fibre of $\int E$ is the discrete set $E$, while its set of total connected components is $E/G$. It follows that a map induced by $h:E\to F$ is fine precisely when $h$ is a bijection, and coarse precisely when $E/G\to F/G$ is a bijection. A bijective equivariant map is already an isomorphism. Hence fine localization of this category changes nothing.

For the coarse localization, write $O(E)=E/G$, and give every set $A$ the trivial $G$-action $T(A)$. The natural equivariant map
\[
 q_E:E\longrightarrow T(O(E))
\]
induces the identity on orbit sets. Thus every $q_E$ is inverted by coarse localization, while $OT=\id_{\mathrm{Set}}$. The induced functors $O$ and $T$ are consequently quasi-inverse equivalences after that localization: $TO$ is naturally isomorphic to the identity through the inverted $q_E$. Under this equivalence the fine-to-coarse functor is exactly $O$.

It remains to verify that the loss of faithfulness persists in the full categories of the proposition, rather than only in the restricted example category. Take $E=G$ with its regular right action. For each $h\in G$, left multiplication $L_h(g)=hg$ is equivariant and defines an endomorphism of $\int G$ over $BG$. Evaluation on the strict fibre followed by $\pi_0$ is a functor
\[
 \Fib_{W_0}(BG)\longrightarrow\mathrm{Set}
\]
which inverts every fine equivalence. It therefore factors through $\HOT_{W_0}^{\mathrm{lc}}(BG)$. Its values on the $L_h$ are distinct set maps, so these endomorphisms remain distinct in the fine localization. On the other hand, $q_G:G\to\{*\}$ is coarse, because both orbit sets are singletons. The equalities $q_GL_h=q_G$ and invertibility of $q_G$ in the coarse localization imply that every $\lambda(L_h)$ is the identity. Thus $\lambda$ is not faithful.\hfill$\square$

This also gives a completely explicit failure of W(4b), even on the perfect objects used in its first weakened variant: $\int G\to BG$ is a total $W_0$-equivalence, whereas its fibre map $G\to\{*\}$ is not. The proposition does not assert that the full coarse category is always a category of sets. The equivalence with sets was proved for the displayed category of discrete fibrations; nonfaithfulness in the full category was established separately by fibre evaluation.

We checked the localizer discussion in \citep{Cisinski2004Minimal} and the locally constant formalism in \citep{Cisinski2009LocallyConstant}. We did not find this exact calculation of the manuscript's $\lambda$ there. The argument is an elementary development of familiar action-groupoid and localization constructions, with no claim of priority.
\end{dcomment}

\begin{dcomment}{exposition}{p180}{Which replacement actually gives a quasi-inverse}
Grothendieck now tests the tools used to draw the preceding diagram. The decisive information is carried by the natural maps into the replacements. A replacement whose value belongs to a desirable subcategory gives an equivalence of localizations only when its comparison map lies in the localizer used on both sides.

For $\Phi$ the comparison $X\to\Phi(X)$ is right-fine on all objects. If $X$ is weakly $W$-fibred, its value is a split locally constant fibration, and the equivalence of fine tests on weakly fibred objects makes the comparison fibrewise. The same reasoning for $\Psi$ uses the left-fine comparison and the proper side. These statements explain why the upper wings remain equivalent after fine localization. They are applications of the localization comparison lemma already discussed in the commentary at manuscript p.~12, rather than a new assertion that localization preserves arbitrary full embeddings. At the split vertices one also uses the free split-fibration adjunction discussed in the commentary at manuscript pp.~53--59; the original unit need not preserve the chosen splitting.

The map $X\to C(X)$ comes with a weaker recorded guarantee: it belongs to $W_S$. Its target is perfect, but perfection of the target alone does not make the comparison fine. The Cartesian square displayed on p.~179 isolates the issue even more sharply. For a total $W$-equivalence $u:X\to Y$, the map $C(u)_s$ is obtained by pulling $u$ back along the perfect map $C(Y)_s\to Y$. To conclude that $C(u)_s\in W$ one needs a suitable preservation theorem for that pullback. Two-out-of-three for the total units only proves $C(u)\in W_S$; it supplies no missing fibre comparison.

There are two different questions here. One asks whether $C$ turns coarse equivalences into fine ones. The other asks whether the inclusion of perfect objects induces an equivalence after fine localization. Failure of the first proposed proof does not establish a negative answer to the second question. Another replacement, or an additional limit argument, could still produce the required fine comparisons. This is exactly what the alternating construction on pp.~192--197 is designed to investigate.

Accordingly, the ``t\'eratologique'' family of localizations is a possibility recorded during the investigation, not a theorem that infinitely many pairwise inequivalent categories have been constructed. The source supplies no family of counterexamples on these pages. Modernly, the alternating replacement settles this particular question: the results explained at p.~172 make its values $W_\omega$-perfect and its units fine on every weakly $W$-fibred input. Every full intermediate category containing $\Parf_\omega S$ therefore has the same fine localization as $\Fib_W S$, by the common replacement and its units. The proof is given at p.~195. This does not identify fine and coarse localization: the connected-components counterexample still applies to $\lambda$.
\end{dcomment}

\begin{dcomment}{exposition}{p181:Ainsi le diagramme $(\ast)$ p. 178 prend la forme}{The fibre axiom closes the comparison}
The force of W(4b) is now visible in one line. For a weakly $W$-fibred object $X$, both $X$ and $C(X)$ are in $\Fib_W S$, and the natural map $i_X:X\to C(X)$ is a total $W$-equivalence. W(4b) makes it a fibrewise $W$-equivalence. This is the missing comparison in the preceding block.

It is useful to check that this gives a functor on the desired localizations, not only good objects. If $u:X\to Y$ is fine, naturality gives
\[
 C(u)i_X=i_Yu.
\]
On each fibre, $i_X$ and $i_Y$ are $W$-equivalences. Two-out-of-three therefore makes $C(u)$ fine. The functor $C$ thus descends to fine localizations. On the full category of weakly fibred objects and on every indicated intermediate full subcategory containing its perfect image, the two composites with inclusion are naturally isomorphic to the identity through the same units. This proves the claimed equivalences.

For $\lambda$ an even shorter argument is available. The two presentations are localizations of the very same category $\Fib_W S$, and W(4b) identifies their two classes of denominators. Choosing the same localization construction on each side makes $\lambda$ literally the identity. Without that choice, the invariant conclusion is its canonical equivalence. Grothendieck's exclamation mark after the displayed equality expresses this formal simplification, not an additional theorem about particular models of localization.

The reduction from four categories to three consequently has an exact hypothesis. The connected-components example above shows that it cannot be asserted for every localizer satisfying only the basic axioms. The next step is to identify the remaining maps as adjoints and to prove that the locally constant category sits fully faithfully inside the two directional theories. Those assertions require the mapping-property calculation on pp.~182--183; they do not follow from drawing a triangle in place of the diamond.
\end{dcomment}

\begin{dcomment}{exposition}{p183}{Local objects absorb the extra denominators}
Grothendieck has reached the point where the three remaining homotopy categories can be compared by their maps, rather than by successive descriptions of their objects. This is the purpose of the boxed formula. Its force is that the source $X$ is arbitrary: local constancy is imposed on the target, where it supplies the extension property needed to invert the additional arrows.

Write $\mathcal H_d(S)=(\Cat/S)[(W_S^{\mathrm d})^{-1}]$ and $\mathcal H_c(S)=(\Cat/S)[W_S^{-1}]$. Under W(4b), the functor $C$ represents a fully faithful right adjoint $i$ to the localization $\rho:\mathcal H_d(S)\to\mathcal H_c(S)$. For $X,Y$ in $\Cat/S$, the adjunction gives
\[
 \Hom_{\mathcal H_d(S)}(X,CY)
 \simeq \Hom_{\mathcal H_c(S)}(X,Y).
\]
This formula is natural in both variables. In particular, it identifies the map induced by localization, after the canonical identification of $CY$ with $Y$ in $\mathcal H_c(S)$; it is not merely an abstract bijection between two sets of the same size. The dual localization gives the other bijection in the text.

To see what changes when $Y$ is already $W$-fibred, use the comparison $i_Y:Y\to CY$. It belongs to $W_S^{\mathrm f}$, and on the category of $W$-fibred objects this is also $W_S^{\mathrm d}$ and $W_S^{\mathrm g}$. Thus postcomposition by $i_Y$ is a bijection on Hom sets in either lateral category. Composing with the displayed adjunction proves the corollary exactly as stated. Notice that no replacement of the source is required. An arbitrary object can map into a local object, and this map already incorporates all the relations that the coarse localization would introduce.

There is a useful intrinsic reformulation. Let $v:A\to B$ be a coarse weak equivalence, and regard it as a map in $\mathcal H_d(S)$. For every $W$-fibred $Y$, precomposition induces a bijection
\[
 v^*:\Hom_{\mathcal H_d(S)}(B,Y)
       \longrightarrow \Hom_{\mathcal H_d(S)}(A,Y).
\]
Indeed the two sides are the corresponding Hom sets in $\mathcal H_c(S)$, where $v$ is invertible. This is the categorical meaning of being local for the additional denominators. It is this property, rather than the mere existence of a map $Y\to CY$, that makes a reflection useful. The distinction was prepared in the commentary at manuscript pp.~64--69; here it becomes an operational formula for computing morphisms.

The hypothesis W(4b) enters at a precise place. Naturality of $X\to CX$ and two-out-of-three show, without W(4b), that $C$ preserves total $W$-equivalences. To conclude that $C(v)$ is also right-fine, one must pass from a total equivalence between fibred objects to equivalences on their fibres. That is exactly what W(4b) supplies. The following pages investigate whether a smaller family of test objects can supply the same passage.
\end{dcomment}

\begin{dcomment}{status}{p184}{The later annotation changes the question}
Grothendieck is asking how much fibre detection is needed to identify the two locally constant theories. The note added to this passage belongs to a later stage of the investigation. In Chapter XVI, Theorem 4.4.13, especially part~$2^\circ$, he proves an equivalence of the relevant fibration conditions when the localizers concerned are stable under sequential colimits. Its conditions $B$, $C$, and $D$ are the later counterparts of the fibre-detection condition, the base-change condition, and the restriction to perfect objects discussed here. The proof uses the alternating replacement of manuscript pages~192--195, together with a factorization argument \citep[\S4.4, especially Corollaries~4.4.10--4.4.12 and Theorem~4.4.13]{Grothendieck1991XVI}.

Thus the apparent weakness of a perfect-object test must be read together with the limit hypothesis. The test cannot be used on an arbitrary weak fibration until a suitable replacement has been constructed and its unit has been proved fine. These are precisely the tasks undertaken later in the chapter. The annotation supplies a conditional resolution in its historical setting. The modern automatic-colimit result discussed after manuscript p.~195 supplies the limit hypothesis for the standing basic localizers, while the restricted fibre-detection test remains necessary. The references to ``XVI~4.9'' in notes~95--96 should be located through the statements: in the current electronic edition the relevant replacement results occur in \S4.4.
\end{dcomment}

\begin{dcomment}{exposition}{p185}{A smaller test family must still contain replacements}
Grothendieck is trying to economize the fibre-detection axiom without losing the reflection. The difficulty is structural. Testing total equivalences only between perfect objects says nothing about the map which first makes a given fibred object perfect, unless that map is already known to be a fibrewise equivalence. This is why the first version of W(4b$'$) does not immediately do the work expected of it.

The issue can be stated without the changing notation. Let $\mathcal F$ be a category of objects over $S$, and let $E\subseteq V$ be two classes of arrows, thought of as fibrewise and total equivalences. Suppose $\mathcal P\subseteq\mathcal F$ is a full subcategory on which $E$ and $V$ agree. This equality alone does not identify $\mathcal F[E^{-1}]$ with $\mathcal F[V^{-1}]$. What is additionally needed is a functor $R:\mathcal F\to\mathcal P$ with a natural map $\eta:\mathrm{id}\to R$ whose components belong to $E$. Once such a functor exists, the conclusion follows, and in fact the argument proves more. For $u:X\to Y$ in $V$, naturality and two-out-of-three imply $R(u)\in V$. Since $RX,RY$ lie in $\mathcal P$, it follows that $R(u)\in E$. In the square
\[
\begin{tikzcd}
 X\arrow[r,"u"]\arrow[d,"\eta_X"'] &Y\arrow[d,"\eta_Y"]\\
 RX\arrow[r,"R(u)"']&RY
\end{tikzcd}
\]
the other three arrows therefore lie in $E$, so $u$ does also. Thus $E=V$ throughout $\mathcal F$, provided these are the actual two-out-of-three classes under consideration; if only invertibility in the localization was assumed, the conclusion is equality of their strong saturations.

In the present application, $\mathcal F=\Fib_W S$, $E=W_S^{\mathrm f}$, and $V=W_S$. The old perfect-object test does not initially provide the required fine unit. By contrast, a test on split fibrations with $W$-equivalent transport does: the functor $\Phi$ takes a $W$-fibred object into that class, and its unit is fibrewise in $W$. The infinite alternating construction below provides another such replacement, now into $\Parf_W S$, assuming the sequential-colimit axiom.

This explains both the frustration and the eventual simplification. There are two different statements to prove: enough good objects exist, and coarse equivalences between these objects are fine. The first is a replacement statement, the second a detection statement. Neither follows merely from writing down a localization. The proposition on manuscript p.~196 combines them; the omitted adjunction check on p.~197 can then be completed by transferring the already established reflector, as we shall do there.
\end{dcomment}

\begin{dcomment}{status}{p186}{What the bifibration question asks}
Grothendieck now asks for a replacement which has both kinds of categorical lifting at once. The question is stronger than the existence of a $W$-perfect replacement. A Grothendieck bifibration has actual Cartesian and coCartesian lifts. Smoothness and properness ask for suitable categories of lifts to be $W$-aspherical. The latter condition can hold without a distinguished universal lift, so an equivalence using $\Parf_W S$ does not by itself answer the question about $\Bifib S$.

One case admits a complete answer with the constructions already at hand. Suppose the base $S$ is a groupoid. Every Grothendieck fibration $p:E\to S$ is then also an opfibration. To prove this, let $u:s\to t$ and $x\in E_s$. Choose a Cartesian lift of $u^{-1}$ with codomain $x$. Its domain lies over $t$, and the lift is an isomorphism because its base arrow is invertible: Cartesian lifting of the inverse gives a two-sided inverse by the uniqueness in the Cartesian universal property. The inverse of this lift is therefore a coCartesian arrow over $u$ with domain $x$. This supplies the required coCartesian lifts for all $u$ and $x$.

Consequently $\Phi(X)\to S$ is a bifibration for every $X\to S$. If $X$ is $W$-fibred, its unit $X\to\Phi(X)$ is in $W_S^{\mathrm f}$. The functor $\Phi$ preserves fine equivalences between such objects: in each fibre it gives a comma functor, and the two fibre-inclusion equivalences identify this with the original fibre map. The same unit, restricted to bifibrations, stays inside that full subcategory. Passing to the fine localizations therefore gives mutually inverse equivalences
\[
 (\Bifib S)[(W_S^{\mathrm f})^{-1}]
 \simeq (\Fib_W S)[(W_S^{\mathrm f})^{-1}]
 \qquad (S\text{ a groupoid}).
\]
No colimit axiom is used. This includes discrete bases, where the assertion is simply the independent comparison on every fibre.

For an arbitrary base, the proof breaks at exactly the attempted inversion of $u$. A fibration classified by $F:S^{\mathrm{op}}\to\Cat$ is a bifibration when its transport functors have the appropriate left adjoints, with the coherent structure coming from the liftings. Knowing only that these transport functors lie in $W$ does not produce such adjoints. The sequential construction below proves the fine-localization comparison for $\Parf_W S$; it does not equip its output with Cartesian and coCartesian lifts, and therefore does not establish the general $\Bifib S$ comparison.

A relevant categorical construction is now available: \citet[Theorem~1.17]{ClarkeSchererZeilberger2025} constructs the free cloven bifibration on any functor $p:X\to S$, with a natural unit $X\to\mathrm{Bif}(p)$. Its universal property uses functors preserving the chosen lifts. Propositions~1.23--1.24 make the unit fully faithful and give it a left inverse on total categories. Those conclusions do not make the unit a fibrewise $W$-equivalence on weakly $W$-fibred inputs, which is the missing assertion for the present application. The antecedent path construction in \citet[\S4, pp.~22--23]{Lamarche2013Path} explicitly distinguishes its own weak equivalences from ordinary nerve weak equivalences. Thus these free constructions cannot be cited as a verified homotopical replacement for this question, even at $W=W_\infty$.

We did not find the precise arbitrary-base, arbitrary-localizer bifibration comparison in the literature checked. This is a limit of the present audit, not an assertion that the problem is open. The modern theory of locally constant coefficients gives a homotopy invariant category of such coefficients; for example, \citet[Theorem~1.3 and Corollary~1.28]{Cisinski2009LocallyConstant} proves invariance under nerve weak equivalences and its detection by locally constant coefficients. That theorem does not assert the strict bifibration presentation asked for here.
\end{dcomment}

\begin{dcomment}{development}{p186}{Bifibrations over posets and over bases with cancellation}
For $W=W_\infty$ the essential surjectivity asked for on this page can be proved for a large class of bases, among them all partially ordered sets. The statement is: every locally constant object of $\mathsf{HOT}(S)$, that is, every diagram of spaces on $S$ whose transition maps are weak equivalences, is the diagram of fibres of a functor $X\to S$ that is at once a Grothendieck fibration and a Grothendieck opfibration. The theorem below was not found in the sources consulted; no priority is claimed.

\textbf{Lemma 1.}
\textit{A Grothendieck fibration is an opfibration if and only if each of its transition functors has a left adjoint. The transition functors of a bifibration are therefore weak equivalences, and a bifibration is proper and smooth.}

This is classical: a cocartesian lift of $\alpha:s\to t$ at $x$ is an initial object of the category of lifts of $\alpha$ with source $x$, which is the comma category of $x$ over the transition functor $\alpha^*:X_t\to X_s$. The converse of the last assertion fails. The four-element poset $x<a,b<c$, mapped to $[1]$ by $x\mapsto0$ and $a,b,c\mapsto1$, is proper and smooth over $[1]$ (the lifting posets $\{a,b,c\}$ and $\{x\}$ have a greatest element) and is a fibration (the cartesian lifts end at $x$); it is not an opfibration, since $x$ has no least element above it in the fibre over $1$.

\textbf{Lemma 2 (patches).}
\textit{Let $Y$ be a simplicial set and $\lambda$ an infinite cardinal with $|Y|\leq\lambda$. Let $K_\lambda(Y)$ be the category of the nonempty weakly contractible simplicial sets of cardinality at most $\lambda$ over $Y$, with all maps over $Y$. The functor $j:\Delta/Y\to K_\lambda(Y)$ which regards a simplex of $Y$ as an object over $Y$ is a weak equivalence. Hence $BK_\lambda(Y)$ is naturally weakly equivalent to $|Y|$.}

\textit{Proof.}
For an object $P\to Y$ of $K_\lambda(Y)$ the comma category $j/P$ is $\Delta/P$, whose classifying space has the weak homotopy type of $P$ and is contractible. Quillen's Theorem~A applies; and $B(\Delta/Y)\to|Y|$ is a natural weak equivalence.\hfill$\square$

\textbf{Lemma 3.}
\textit{Let $f:Y\to Z$ be a trivial Kan fibration of simplicial sets of cardinality at most $\lambda$. Then composition with $f$, $f_!:K_\lambda(Y)\to K_\lambda(Z)$, has the right adjoint $f^*$, pullback along $f$.}

\textit{Proof.}
For $Q\to Z$ in $K_\lambda(Z)$ the projection $Q\times_ZY\to Q$ is a trivial Kan fibration, so $Q\times_ZY$ is nonempty and weakly contractible, of cardinality at most $\lambda$. The adjunction is that of the slice categories.\hfill$\square$

\textbf{Theorem.}\ours
\textit{Let $S$ be a small category in which every arrow is a monomorphism, or every arrow is an epimorphism; for instance a partially ordered set. Then every locally constant object of $\mathsf{HOT}(S)$ is represented by the diagram of fibres of a bifibration $X\to S$ with $X$ small.}

\textit{Proof.}
Suppose first that every arrow is a monomorphism, and represent the object by a functor $F:S^{\mathrm{op}}\to\mathbf{sSet}$ whose transition maps are weak equivalences. Let $F\to RF$ be a fibrant replacement for the injective model structure on simplicial presheaves on $S$, in which the weak equivalences are objectwise and the cofibrations are the objectwise monomorphisms. We claim that for every arrow $\alpha:s\to t$ the map $RF(\alpha):RF_t\to RF_s$ is a trivial Kan fibration. Let $h_s\to h_t$ be the map of representable presheaves induced by $\alpha$; it is objectwise injective because $\alpha$ is a monomorphism. A lifting problem of a horn inclusion $\Lambda^n_k\subset\Delta^n$ against $RF(\alpha)$ is, by the Yoneda lemma, a lifting problem of
\[
 (h_t\times\Lambda^n_k)\amalg_{h_s\times\Lambda^n_k}(h_s\times\Delta^n)\longrightarrow h_t\times\Delta^n
\]
against $RF\to*$. At each object this map is a disjoint union of identities of $\Delta^n$ and of horn inclusions; it is an objectwise trivial cofibration, and $RF$ is injectively fibrant. So $RF(\alpha)$ is a Kan fibration; it is a weak equivalence because $F(\alpha)$ is one and $F\to RF$ is an objectwise weak equivalence.

Choose $\lambda$ bounding the values of $RF$. By Lemma~3 the functor $s\mapsto K_\lambda(RF_s)$, with the transition functors $RF(\alpha)_!$, takes values in functors having right adjoints; so $s\mapsto K_\lambda(RF_s)^{\mathrm{op}}$ is a functor on $S^{\mathrm{op}}$ whose transition functors have left adjoints. Its Grothendieck construction $X\to S$ is a fibration, and by Lemma~1 a bifibration. By Lemma~2, and because a category and its opposite have the same classifying space naturally, its diagram of fibres is naturally weakly equivalent to $RF$, hence to $F$.

If every arrow is an epimorphism, the covariant representables give objectwise injections $S(t,-)\to S(s,-)$, and the same argument applies to a covariant model of the local system, with Lemma~3 used on the patch categories themselves.\hfill$\square$

Two further cases hold over every small base $S$, by simpler arguments which are not reproduced here: when the values of the local system are 1-types (take fundamental groupoids objectwise; equivalences of groupoids have adjoints), and when $BS$ is a 1-type (pull back from the groupoid obtained by inverting all arrows of $S$; this contains the case of a groupoid treated above, and the bases with an initial or a terminal object).

\textbf{What remains.}
For a general base the question is undecided. The first test case is a monoid $M$ with five elements, the identity and a rectangular band $\{e_{ij}\}_{i,j\in\{1,2\}}$ with $e_{ij}e_{kl}=e_{il}$, whose classifying space has the homotopy type of $S^2$, with the local system of the path fibration (fibre $\Omega S^2$). The method above cannot reach it. A strict action of a monoid on a simplicial set in which an idempotent $e$ acts by a trivial Kan fibration has $e$ acting as the identity, because a surjective idempotent map is the identity; so a strict diagram with trivial fibrations as transitions on a monoid generated by idempotents is constant, and a bifibration for the test case must come from a pseudofunctor that is not strict. The second half of the question of this page, whether the comparison of localizations is an equivalence, is not addressed by the theorem.
\end{dcomment}

\begin{dcomment}{examples}{p190}{A branching base defeats one-sided completion}
Grothendieck needs a counterexample which separates transport in one direction from transport in the other. A base with a terminal object hides this distinction: an opfibration over such a base has the homotopy type of its terminal fibre. The branching poset $S=(0\to a,\,0\to b)$ removes that shortcut.

Take a pointed category $(T,t)$, and form the split opfibration classified by
\[
 e\xrightarrow{t}T,\qquad e\xrightarrow{\mathrm{id}}e.
\]
Its total category $X$ has the objects of $T$, together with objects $\xi_0$ over $0$ and $\xi_b$ over $b$. The additional arrows are a single arrow $\xi_0\to\xi_b$ and, for every $z\in T$, arrows $\xi_0\to z$ indexed by $\Hom_T(t,z)$. Composition with arrows of $T$ is the evident one. There are no arrows from $T$ to $\xi_0$ or $\xi_b$, and none from $\xi_b$ to the other objects.

Let $X''$ be the full subcategory obtained by deleting $\xi_b$. There is a retraction $r:X\to X''$ which sends $\xi_b$ to $\xi_0$ and the arrow $\xi_0\to\xi_b$ to an identity. If $j:X''\hookrightarrow X$ denotes the inclusion, then $rj=\mathrm{id}$ and there is a natural transformation $jr\Rightarrow\mathrm{id}_X$, equal to the identity on $X''$ and to $\xi_0\to\xi_b$ at $\xi_b$. Naturality needs checking only on the new arrow, where both composites are that arrow. Thus $j$ is a categorical homotopy equivalence and belongs to every localizer in use here.

There is a second retraction $r_a:X''\to T$, sending $\xi_0$ to $t$ and fixing $T$. In fact $r_a$ is left adjoint to the inclusion $j_a:T\hookrightarrow X''$: the only nontrivial adjunction formula is
\[
 \Hom_T(r_a\xi_0,z)=\Hom_T(t,z)
              =\Hom_{X''}(\xi_0,j_a z).
\]
Hence $T\to X''\to X$ belongs to $W$. The fibre over $b$ is nevertheless $e$. If $T\to e$ is not in $W$, the map $X_b\to X$ is not in $W$, by two-out-of-three. Since $\Phi(X)_0\simeq X$ and $\Phi(X)_b\simeq e$, transport along $0\to b$ in the completed fibration is not a $W$-equivalence. This proves that $\Phi(X)$ need not even be $W$-fibred, although $X\to S$ is an opfibration and hence proper.

For the ordinary weak-homotopy localizer, a completely finite choice is $T=\{t,t'\}$ discrete. Then $X$ is the disjoint union of the three-object poset $\xi_b\leftarrow\xi_0\to t$ and the isolated object $t'$. Its nerve has two contractible components, whereas $X_b$ has one. No homology calculation or pushout theorem is required. For a general $W$, the argument requires a pointed $T$ which is not $W$-aspherical; the nondegeneracy assumption $W\ne\mathrm{Fl}(\Cat)$ alone should not be substituted for that requirement in this particular example.

The parenthetical pushout question has a positive answer under the suggested symmetric hypotheses. If $A\hookrightarrow B$ and $A\hookrightarrow C$ are both full sieves, their pushout $P$ contains no arrows from $B\setminus A$ to $C\setminus A$, or conversely. Every composable chain lies in one side, so $NP=NB\cup_{NA}NC$. The nerve inclusions are monomorphisms, and this union is the homotopy pushout of the two inclusions. Equivalently one can check the canonical functor from the total category of $B\leftarrow A\to C$ by comma categories: for $p\in B\setminus A$, retract the copy of $A/p$ lying on the $C$-side onto the middle copy. The resulting category has terminal object $p$ on the $B$-side. The case $p\in C\setminus A$ is symmetric, and for $p\in A$ the comma category contracts to the three copies of $p$ joined by a span. Theorem A therefore proves the comparison is in $W$. For two cosieves, apply the dual argument. The special example above is simpler still, since the explicit retractions determine the homotopy type for every allowed $W$.
\end{dcomment}

\begin{dcomment}{development}{p192}{The obstruction persists on any base with two maximal branches}
Grothendieck strengthens the preceding counterexample: over the three-object branching poset, a single application of $\Phi$ cannot be proper unless the original category is empty. The obstruction lies in the absence of arrows between the two maximal branches. We can isolate that reason and extend the statement without introducing additional homotopical machinery.

\textbf{Branching obstruction.}\ours\enspace
Let $S$ be a poset with a least element $0$ and two distinct maximal elements $a,b$. Let $W$ be a localizer satisfying the standing axioms and suppose that $\varnothing\to e$ does not belong to $W$. For every $X\to S$, if $\Phi_S(X)\to S$ is $W$-proper, then $X=\varnothing$.

\emph{Proof.} The fibre diagram of $\Phi_S(X)$ is $s\mapsto s\backslash X$. Since $0$ is least, $0\backslash X$ identifies with $X$. Since $a$ and $b$ are maximal, $a\backslash X=X_a$ and $b\backslash X=X_b$. The two transport functors to the fibre at $0$ are therefore the inclusions $X_a\hookrightarrow X$ and $X_b\hookrightarrow X$.

Properness forces each of these inclusions to be coaspherical. Here is the precise test being used. Pull the fibration back along the arrow $0\to a$. For an object $x$ of its fibre over $0$, the category of arrows from $x$ to the fibre over $a$ is the comma category $x\backslash X_a$. Properness requires it to be $W$-aspherical. This is the elementary lifting-category criterion for properness, applied to this one arrow; the same statement holds for $b$.

If $x\in X$ existed, the category $x\backslash X_a$ would be $W$-aspherical and hence nonempty. Choose one of its objects, an arrow $x\to x_a$ with $x_a\in X_a$. Properness along $0\to b$, applied now to the object $x_a$ of $X$, would require $x_a\backslash X_b$ to be $W$-aspherical. But this category is empty: an arrow $x_a\to x_b$ would project to an arrow $a\to b$, impossible for distinct maximal elements of a poset. This contradicts the hypothesis on the empty category. Thus $X$ has no objects. Conversely, the empty category over $S$ is proper, since all the lifting tests with a specified source object are vacuous.\hfill$\square$

For the standing localizers, the nondegeneracy condition in the text implies the displayed hypothesis on $\varnothing\to e$. Indeed, if $\varnothing\to e$ were in $W$, product stability would put $\varnothing\to A$ in $W$ for every category $A$. Two-out-of-three would then put every functor in $W$.

This proof also fills the final step in the boxed assertion on manuscript p.~192. Showing that $X_a=X_b=\varnothing$ is not yet, by itself, a proof that $X_0=\varnothing$ for an arbitrary $X\to S$. One must use coasphericity once more: $\varnothing\to X$ cannot be coaspherical when $X$ has an object. The argument above incorporates this step directly. The preceding contravariant diagram prints $X_0$ at its centre; its mathematical value is $X=0\backslash X$, as both the following inclusions and the proof require. That printed slip is retained in the source and recorded in the correction ledger.

We checked the lifting criteria in \citet[Proposition~4.3, Theorem~4.13, and Example~4.17]{Maltsiniotis2005Asphericite}; we did not find this particular branching formulation there. The proposition is an elementary extension of the present manuscript argument, with no claim of priority.
\end{dcomment}

\begin{dcomment}{exposition}{p194}{Why alternating completion works at the limit}
The failure of one-sided completion suggests Grothendieck's next move. Instead of asking $\Phi$ to preserve properness, he restores properness with $\Psi$, then restores smoothness with $\Phi$, and continues. No finite stage has to possess both properties. At the limit, the odd and even stages are cofinal and can establish the two properties separately.

Put $X_0=X$, $X_1=\Phi X$, $X_2=\Psi\Phi X$, and continue alternately, with the canonical unit at each step. Let $R X=\varinjlim_n X_n$. W(6 bis) means that an objectwise $W$-equivalence between two sequential diagrams induces a $W$-equivalence on their ordinary colimits in $\Cat$. Its quantifier is sequential. The unqualified words ``limite inductive filtrante'' in the next paragraph must be read with this restriction; an arbitrary filtered indexing category is not covered by the axiom as written.

The passage from this axiom to smoothness and properness of $RX\to S$ can be made explicit by the lifting-category criterion. For a functor $p:E\to S$, a base arrow $u:s\to t$, and $y\in E_t$, let $L_p(u,y)$ have objects arrows $v:x\to y$ with $p(v)=u$ and morphisms commuting triangles whose arrow between sources lies over $\mathrm{id}_s$. Smoothness is equivalent to $L_p(u,y)$ being $W$-aspherical for every $u,y$. Properness is the dual condition on arrows with a specified source. These are the familiar fibre categories of $E/y\to S/t$ and their duals.

Suppose first that $E=\varinjlim_n E_n$ over the fixed base $S$, with every $E_n\to S$ smooth. An object $y$ of $E$ is represented by $y_n\in E_n$ at some stage. After restricting to that tail of the sequence, there is a canonical isomorphism
\[
 L_p(u,y)\simeq\varinjlim_{m\ge n}L_{p_m}(u,y_m).
\]
To verify it, represent an arrow to $y$ at a sufficiently late stage; its image in $S$ is already the required fixed arrow $u$. A morphism of the lifting category is represented by one further arrow and one equality of composites, and that equality holds at a later stage. The same argument detects equality of two such morphisms. Thus the displayed functor is bijective on objects and arrows after taking the colimit. Every stage maps to $e$ by a member of $W$, so W(6 bis) makes the colimit $W$-aspherical. This proves sequential stability of smoothness. Reversing arrows proves stability of properness.

Now the odd subsequence $X_1,X_3,\ldots$ consists of smooth categories over $S$, and the even subsequence $X_2,X_4,\ldots$ consists of proper ones. Both have colimit $RX$, because each is cofinal in the original sequence. Hence $RX$ is $W$-perfect. Every transition $X_n\to X_{n+1}$ belongs to $W_S$, so every composite $X\to X_n$ does. Apply W(6 bis) to the natural transformation from the constant sequence at $X$ to $(X_n)$: its colimit is $X\to RX$, proving that this map lies in $W_S$ for arbitrary $X$.

When $X$ is $W$-fibred, all stages remain $W$-fibred and the transitions are fibrewise $W$-equivalences. Fibres over a fixed $s$ commute with these colimits: their objects and arrows are defined by fixed equations in $S$, and any equality in a filtered colimit is witnessed at a later stage. Applying W(6 bis) in each fibre yields
\[
 X_s\longrightarrow (RX)_s\quad\text{in }W
 \qquad\text{for every }s\in S.
\]
Thus the same construction supplies both kinds of unit advertised in the text: a total equivalence for arbitrary objects and a fine equivalence for fibred objects. It is the latter statement that was missing for the earlier proposed perfect-object replacement. Nothing in this argument asserts that $RX\to S$ is a Grothendieck bifibration: a sequential colimit of lifting categories known to be aspherical does not supply actual universal lifts.

For $W_\infty$ the two assertions hold without the axiom W(6 bis) as a separate hypothesis: the lifting categories of a filtered colimit are the filtered colimits of the lifting categories, and filtered colimits of weakly contractible simplicial sets are weakly contractible. The limit is proper and smooth; it need not be a bifibration, even when every stage is one. Let $X_n$ be the chain $0<1<\dots<n<t$ over $[1]$, with $t$ over $1$ and the rest over $0$. Each $X_n\to[1]$ is a bifibration (the cartesian lift of $0\to1$ ends at $n$). The colimit of the inclusions $X_n\subset X_{n+1}$ has the fibre $\mathbf N$ over $0$, and a cartesian lift with target $t$ would be a greatest natural number. This is why the alternating construction $\Phi_\infty$ yields perfect objects and not bifibrations.
\end{dcomment}

\begin{dcomment}{exposition}{p195}{Fine replacement gives an equivalence of localizations}
Grothendieck can now return to the question on manuscript p.~186 with an actual construction in hand. The essential point is stronger than the existence of a perfect object with the right total homotopy type: for a fibred input, the comparison is an equivalence in every fibre. This is what allows the replacement to control morphisms as well as objects.

Today W(6 bis) is automatic: every basic localizer is closed under filtered colimits \citep[Corollary~4.2.22]{Cisinski2006}. This applies to the standing relative W$(1,2,3)$, without accessibility; we do not extend it here to arbitrary weak basic localizers.

Let $R=\Phi_\infty$, let $\mathcal F=\Fib_W S$, and let $\mathcal P=\Parf_W S$. Under W(6 bis), the unit $\eta_X:X\to RX$ is fine for every $X\in\mathcal F$. The functor $R$ preserves fine equivalences on $\mathcal F$. Indeed, if $u:X\to Y$ is fine, then in each fibre the equality
\[
 (Ru)_s(\eta_X)_s=(\eta_Y)_s u_s
\]
and two-out-of-three show that $(Ru)_s$ lies in $W$. Thus both the inclusion $J:\mathcal P\hookrightarrow\mathcal F$ and $R$ induce functors after fine localization. The units give natural isomorphisms
\[
 \mathrm{id}\xrightarrow{\sim}\overline J\,\overline R
 \quad\text{on }\mathcal F[(W_S^{\mathrm f})^{-1}],
 \qquad
 \mathrm{id}\xrightarrow{\sim}\overline R\,\overline J
 \quad\text{on }\mathcal P[(W_S^{\mathrm f})^{-1}].
\]
For the second assertion, both endpoints of the unit belong to $\mathcal P$, and the subcategory is full. These two natural isomorphisms prove that the induced functors are quasi-inverse equivalences. In particular, full faithfulness is part of the result; there is no separate assumption about roofs or a calculus of fractions.

The same proof works with any full intermediate subcategory containing the image of $R$. The more precise claim in footnote~94 follows because the construction already lands in $\Parf_{W''_\omega}S$. One must keep the two minimal classes distinct at this stage of the manuscript. The inclusion $W_\omega\subseteq W''_\omega$ does not, on its own, put the output in $\Parf_\omega S$. It puts it in the class named in the footnote.

The minimality comparison above does, however, allow us to settle the stronger annotation in footnote~89. It identifies $W_\omega$ with the nerve weak-homotopy equivalences $W_\infty$. This class is stable under sequential colimits. Indeed, the nerve commutes with filtered colimits because each finite ordinal category $[n]$ is finitely presented. Weak equivalences of simplicial sets are stable under such colimits: apply $\operatorname{Ex}^{\infty}$, which commutes with them since every subdivided simplex is finite; a filtered colimit of Kan complexes is Kan, and maps and homotopies from finite simplicial spheres are witnessed at finite stages, giving the required isomorphisms on all homotopy groups and on components. Thus $W_\infty$ satisfies W(6 bis). Minimality gives
\[
 W_\omega\subseteq W''_\omega\subseteq W_\infty=W_\omega.
\]
Consequently the replacement can indeed be taken in $\Parf_\omega S$, with the fine unit just proved. This conclusion uses the modern identification of the minimal class; it was not a consequence of the single inclusion written on manuscript p.~195.

The original restricted W(4b$'$) now suffices as well. Given a total $W$-equivalence $u:X\to Y$ between weakly fibred objects, naturality of the unit makes $Ru$ a total $W$-equivalence between $\Parf_\omega$-objects. The restricted hypothesis makes $Ru$ fine, and the two fine units then make $u$ fine by two-out-of-three in each fibre. This proves full W(4b), and therefore the fine-to-coarse comparison. It explains precisely how the historical retraction acquires a positive answer once the later colimit and minimality theorems are available.

This replacement equivalence still does not identify the fine and coarse localizations. It allows either of them to be computed on a smaller category; it does not say that the two classes of denominators agree there. That remaining assertion is exactly the restricted W(4b) hypothesis in the following proposition. Grothendieck's distinction between replacing objects and comparing equivalences is therefore preserved even after the successful infinite construction.
\end{dcomment}

\begin{dcomment}{exposition}{p197}{Completing the omitted compatibility of the two units}
Grothendieck has proved enough to compare the localizations, but deliberately leaves the triangle check for $\Phi_\infty$ unfinished. There is a short way to complete it which uses the fine replacement already constructed. It also explains the later annotation that the proposed weakened hypotheses recover the full fibre-detection axiom.

First consider any $S$-map $u:X\to Y$ between $W$-fibred categories, with $u\in W$. Under condition~$1^o$, put $R=\Phi$; under condition~$2^o$, put $R=\Phi_\infty$. In either case $\eta_X$ and $\eta_Y$ are fine equivalences. Naturality and two-out-of-three show $Ru\in W$. In the first case $RX,RY$ are split fibrations with $W$-equivalent transports; in the second they lie in $\Parf_{W''_\omega}S$. The restricted fibre-detection hypothesis therefore gives $Ru\in W_S^{\mathrm f}$. Applying two-out-of-three in each fibre to
\[
 (Ru)_s(\eta_X)_s=(\eta_Y)_s u_s
\]
gives $u_s\in W$. Thus each displayed hypothesis implies W(4b) itself. In the first case, it is enough to impose the test only on Cartesian maps, since $\Phi(u)$ is Cartesian. This is the replacement argument formalized later in \citet[Lemma~4.4.9 and Corollary~4.4.10]{Grothendieck1991XVI}; the calculation here uses only the constructions already present in Chapter~VI.

Under condition~$1^o$, the established reflector $C$ already proves the proposition. We now work under condition~$2^o$, retaining W(6 bis), and compare $R=\Phi_\infty$ with $C$. Write $\alpha_X:X\to CX$ for the unit of $C$. Work in $\mathcal H_d(S)$, where the established reflection gives $C\simeq i\rho$. Let $R=\Phi_\infty$ and retain $\eta_X:X\to RX$ for its unit. The object $RX$ is $W$-fibred, so $\alpha_{RX}$ is an isomorphism in $\mathcal H_d(S)$. Also $\eta_X\in W_S$, hence $C(\eta_X)$ is right-fine under W(4b). We therefore have a natural isomorphism
\[
 \theta_X=(\alpha_{RX})^{-1}C(\eta_X):CX\xrightarrow{\sim}RX.
\]
Naturality of $\alpha$ gives the exact identity
\[
 C(\eta_X)\alpha_X=\alpha_{RX}\eta_X,
 \qquad\text{hence}\qquad \theta_X\alpha_X=\eta_X.
\]
Thus $R$ represents the same reflection, with its given unit. This is stronger than an unspecified isomorphism between the two endofunctors: it identifies the units, which are what determine the adjunction.

It remains to extract the equality Grothendieck asks for. Since $R^2X$ is a local object, the universal property of the reflection says that precomposition by $\eta_X$ is a bijection
\[
 \Hom_{\mathcal H_d(S)}(RX,R^2X)
       \xrightarrow{\sim}
 \Hom_{\mathcal H_d(S)}(X,R^2X).
\]
The two maps $\eta_{RX}$ and $R(\eta_X)$ have the same image under this bijection, because naturality of $\eta$ gives
\[
 \eta_{RX}\eta_X=R(\eta_X)\eta_X.
\]
Therefore $\eta_{RX}=R(\eta_X)$ in $\mathcal H_d(S)$, as required. Both maps are isomorphisms: the first because $RX$ is fibred, the second because $R$ sends coarse equivalences to fine ones after the fibre-detection argument. Inverting this common map supplies the multiplication of the idempotent reflector, and the triangle identities follow.

The missing verification is consequently a genuine omission in the written proof, but no additional axiom is needed to repair it. The repair is not the generally invalid cancellation of $\eta_X$ from an equality of composites; $\eta_X$ need not be invertible in the lateral category. Cancellation is justified here by its proven universal property against the local target $R^2X$.
\end{dcomment}

\begin{dcomment}{exposition}{p199}{Restriction depends on the chosen equivalences}
Grothendieck's three inverse-image functors are built from the same categorical pullback, but they are defined on different useful representatives. This is why the functorial diagram is more informative than an undifferentiated assertion that pullback preserves weak equivalences.

For a smooth representative $X\to S$, its pullback along $f:S'\to S$ is again smooth, and fibrewise equivalences remain fibrewise equivalences. Since fine and right-fine equivalences agree on smooth objects, this constructs $f^*$ on $\HOT_W$. The dual argument uses proper representatives for $\HOT_W^o$. On $\Fib_W S$, the same pullback preserves the class of objects and the fine equivalences, and therefore constructs $f^*$ on $\HOT_W^{\mathrm{lc}}$. In each case the composition and identity constraints are those of categorical pullback, passed through localization. Their coherence is inherited from the coherent pullback identifications before localization.

For the coarse localization, this argument is unavailable until total equivalences are controlled. If $f$ is universally in $W$, every projection $X\times_S S'\to X$ belongs to $W$. Given $u:X\to Y$, the natural square with these projections proves
\[
 u\in W\quad\Longleftrightarrow\quad
 u\times_S S'\in W
\]
by two-out-of-three. This proves both preservation and reflection of total equivalences; no good-representative restriction is needed. Under W(4a), the same preservation statement holds when $f$ is $W$-fibred, by the defining universal preservation property of that axiom. It then computes the coarse restriction directly on every $X\to S$.

There is a second, differently placed improvement. If $p:S'\to S$ is proper, pullback preserves right-fine equivalences on all of $\Cat/S$. The displayed formula for $p^*$ on an arbitrary representative is therefore legitimate. The condition is on the map changing the base, whereas the first construction required smoothness of the object being pulled back. Either condition can make the ordinary pullback compute the desired derived inverse image, but the two arguments test different arrows.

These distinctions also explain the compatibility claims. Restriction on fine representatives commutes with the inclusions of locally constant objects into the lateral theories. Commutation with the coarse localization requires preservation of the coarse denominators, and therefore the extra hypotheses just discussed. A diagram of functors cannot be promoted from one assertion to the other merely because the object formulas look identical.
\end{dcomment}

\begin{dcomment}{exposition}{p200}{Multiplication by a category over the base}
Grothendieck introduces $\mathrm{mult}_{X/S}$ because an arbitrary category over the base still acts on relative homotopy types, even when its strict fibre product does not compute the internal product. The action remembers the map $X\to S$ before replacement.

Let $p:X\to S$. For a smooth representative $Y\to S$, pullback computes $p^*[Y]$ as $[X\times_S Y]$ over $X$. Postcomposition then gives
\[
 \mathrm{mult}_{X/S}=p_!p^*.
\]
This identity already explains why the functor is well defined for arbitrary $X$: the right-fine equivalences are respected in the two stages where Grothendieck factors the construction. The repeated $W_S^{\mathrm d}$ at the start of printed p.~124 is a source slip; the pullback stage has source class $W_S^{\mathrm d}$ and target class $W_X^{\mathrm d}$.

By contrast, the internal product is obtained by replacing both factors by smooth categories:
\[
 [X]\times_S[Y]=[\Phi X\times_S\Phi Y].
\]
The map $X\to\Phi X$ induces a comparison from the action to multiplication by $[X]$. For smooth $Y$, it is represented by $X\times_S Y\to\Phi X\times_S Y$. Its failure to be right-fine in general is the counterexample already analyzed on manuscript pp.~153--156; there is no need to repeat it here. What changes in the present discussion is the proposed hypothesis that $X$ is $W$-fibred.

The argument given on pp.~201--202 correctly proves the action formula when $X$ is fibred. Indeed the unit $X\to\Phi X$ is fine; after arbitrary base change to $Y$ it is fine between $W$-fibred objects over $Y$, hence right-fine there. Forgetting the intermediate base preserves right-fine equivalences. This proves that the action of $Y/S$ on $[X]$ can be computed by the strict fibre product even for arbitrary $Y$.

There remains an apparent asymmetry: identifying that value with $[X]\times_S[Y]$ also requires replacing $Y$. Grothendieck proves this when $X$ is proper and suspects that the extra condition is necessary. The next development supplies the missing comparison for a $W$-fibred $X$ by using the transports in its split-fibration replacement. The argument will identify the canonical comparison map itself, so it also gives the natural compatibility needed for products and inverse images.
\end{dcomment}

\begin{dcomment}{development}{p202}{A fibred factor suffices for the product formula}
Grothendieck's desired formula on manuscript p.~201 is stronger than the proof immediately following it. The remaining difficulty can be resolved by calculating over a comma category: the arrow from its fixed starting object trivializes the transport of a locally constant fibration. This gives the unrestricted second factor that the manuscript leaves undecided.

\textbf{Product comparison for a fibred factor.}\ours\enspace
Assume W$(1,2,3)$. If $X\to S$ is $W$-fibred in the weak, universally stable sense used in this chapter, then for every $Y\to S$ the canonical map
\[
 [X\times_S Y]\longrightarrow [X]\times_S[Y]
 \quad\text{in }\HOT_W(S)
\]
is an isomorphism. No properness assumption on $X$ or smoothness assumption on $Y$ is required.

\emph{Proof.} We first treat a split Grothendieck fibration $X=\int_S F$ with $F:S^{\mathrm{op}}\to\Cat$ and $F(u)\in W$ for every arrow $u$ of $S$. Fix $s\in S$ and put $B=s\backslash Y$. An object of $B$ is $(y,a)$, with $a:s\to q(y)$, where $q:Y\to S$ is the structure map. The projection
\[
 E=s\backslash(X\times_S Y)\longrightarrow B
\]
is the split fibration classified by $(y,a)\mapsto F(qy)$. A morphism $(y,a)\to(y',a')$ is an arrow $h:y\to y'$ with $qh\,a=a'$, and its transport is $F(qh)$.

The maps $F(a):F(qy)\to F(s)$ form a natural transformation to the constant diagram at $F(s)$, because
\[
 F(a)F(qh)=F(qh\,a)=F(a').
\]
They therefore define a Cartesian functor
\[
 T:E\longrightarrow F(s)\times B,
 \qquad (x,y,a)\longmapsto (F(a)x,(y,a)).
\]
Each fibre map of $T$ belongs to $W$. Hence $T\in W$: a fibrewise $W$-equivalence between Grothendieck fibrations is right-fine, as seen by the fibre-to-comma adjunctions, and the relative localization axiom sends right-fine equivalences to total ones. This implication uses the fibrewise-to-total direction; it does not invoke W(4b).

We check that this proves the canonical comparison, not just an equivalence between its endpoints. Since $X$ is smooth, $[X]\times_S[Y]$ is represented by $X\times_S\Phi Y$. Let
\[
 k_s:s\backslash(X\times_S Y)
       \longrightarrow s\backslash(X\times_S\Phi Y)
\]
be induced by the unit of $Y$. The target is a comma category of a split fibration, so pullback to its fibre at $s$ defines a functor
\[
 r_s:s\backslash(X\times_S\Phi Y)
          \longrightarrow F(s)\times(s\backslash Y).
\]
The inclusion of that fibre is left adjoint to $r_s$: an arrow from a fibre object to an object over $t$, lying over $c:s\to t$, is uniquely an arrow in the fibre to its Cartesian pullback. Thus $r_s\in W$. On an object $(x,y,a)$ of the source of $k_s$, the unit of $Y$ introduces the identity arrow at $q(y)$, and $r_s$ pulls back along $a$. Consequently $r_s k_s=T$, on objects and on morphisms. Since $r_s,T\in W$, two-out-of-three gives $k_s\in W$. This holds for every $s$, proving
\[
 X\times_S Y\longrightarrow X\times_S\Phi Y
       \quad\text{belongs to }W_S^{\mathrm d}.
\]

For a general $W$-fibred $X$, its replacement $\Phi X$ is a split fibration with $W$-equivalent transports. The unit $X\to\Phi X$ is fine. After pullback to any $Y$, it is fine between $W$-fibred objects over $Y$, because this class of objects is stable under arbitrary base change. Fine and right-fine equivalences agree on that class, and forgetting $Y$ preserves right-fine equivalences. Hence
\[
 X\times_S Y\longrightarrow\Phi X\times_S Y
       \quad\text{belongs to }W_S^{\mathrm d}.
\]
Apply the split case to $\Phi X$. Composing the two maps gives a right-fine equivalence from $X\times_S Y$ to $\Phi X\times_S\Phi Y$, whose image is exactly the canonical product comparison. This proves the assertion.\hfill$\square$

The universal stability used in the last paragraph is part of the weak-fibration notion, not a hidden strengthening of it; see \citet[Chapter~XII, Part~III, Theorem~3, manuscript pp.~51--52]{Grothendieck1990XIIa}. The construction also shows naturality in both $X$ and $Y$, since each map was induced by the functorial comma replacement or by its Cartesian transport.

In particular, the proposed general formula $(*)$ on manuscript p.~201 is valid under its stated fibred hypothesis. The suspicion at the end of p.~202 that properness of $X$ or smoothness of $Y$ is necessary is not borne out. We checked the comma-category and base-change treatments in \citet[\S\S4--5]{Maltsiniotis2005Asphericite} and the locally constant framework in \citet{Cisinski2009LocallyConstant}; we did not locate this exact formulation there. The proof is included as a continuation of the manuscript's argument, without a claim that the result is new.
\end{dcomment}

\begin{dcomment}{later}{p204}{The final adjunction and the surviving base-change fragment}
Grothendieck closes the revision by separating the constructions that need only the three standing axioms from the stronger locally constant comparison. Postcomposition $f_!$ respects the three relative equivalence classes, and the earlier unit and counit prove $f_!\dashv f^*$ for $\HOT_W$ without W(4b). The locally constant adjunction requires the additional comparison that identifies the fine theory with the coarse localization.

For $f\in W$ under W(4b), the locally constant unit and counit are isomorphisms; hence $f_!^{\mathrm{lc}}$ and $f^*_{\mathrm{lc}}$ are inverse equivalences. This assertion concerns locally constant coefficients. It does not say that restriction on all coherent diagrams is an equivalence whenever the indexing categories have the same homotopy type. For instance, the terminal category and the walking arrow have equivalent nerves, whereas the latter still indexes a genuine coherent morphism. The distinction is already built into the separate lateral and locally constant categories.

The constant-coefficient claim has an elementary calculation. If $p:S\to e$ and $T\in\Cat$, the constant coefficient represented by $T$ is represented over $S$ by $S\times T$. Its homotopy colimit is therefore represented by $S\times T$ as an absolute category. A functor $f:S'\to S$ acts by $f\times\mathrm{id}_T$. Product stability gives the implication $f\in W\Rightarrow f\times T\in W$. Conversely, testing $T=e$ detects whether $[f]$ is invertible in $\Cat[W^{-1}]$. This explains the appearance of the strong saturation $\overline W$ in the final statement, rather than an assertion that every chosen class $W$ was strongly saturated from the outset.

The last paragraph of the source displays a square and the hypotheses ``$f$ smooth or $g$ proper'', then stops. Its missing continuation is a source lacuna, not an unproved conclusion which we may attribute to Grothendieck. If the square is Cartesian, the expected statement can be read against the base-change result proved earlier in the chapter: the canonical mate
\[
 f'_!q^*\longrightarrow g^*f_!
\]
is invertible for such a square. The surviving fragment does not itself repeat the Cartesian hypothesis; it is necessary in this reconstruction. In the smooth-$f$ case, represent the coefficient by $Z\to X$ smooth over $X$. Its composite to $S$ is smooth, so both sides are represented by the same category $Z\times_S S'$ over $S'$. In the proper-$g$ case, ordinary pullback along $g$ computes inverse image for every representative, and its pullback $q$ is proper as well; the same equality of representatives proves the formula. The map is the standard mate built from the unit and counit, so the identification agrees with the canonical base-change comparison rather than merely producing some isomorphism between the functors.

This is the categorical form of the base-change mechanism isolated in \citet[Proposition~5.3, Theorem~5.4, and Remark~5.5]{Maltsiniotis2005Asphericite}, with the opposite-category convention adjusted to the present $s\backslash X$ model. The surviving text is left unchanged at its stopping point; the preceding calculation belongs to the commentary.
\end{dcomment}
